% !TEX program = xelatex
% 任务规划算法体系 - 案例式教学版
% 编译方式: xelatex main_casebased.tex (运行两次以生成目录)

\documentclass[a4paper,12pt,openany]{ctexbook}

% ============ 基础宏包 ============
\usepackage[top=2.5cm,bottom=2.5cm,left=2.5cm,right=2.5cm]{geometry}
\usepackage{amsmath,amssymb,amsthm}
\usepackage{graphicx}
\usepackage{booktabs}
\usepackage{longtable}
\usepackage{tabularx}
\usepackage{multirow}
\usepackage{enumitem}
\usepackage{tikz}
\usetikzlibrary{shapes,arrows,positioning,calc,fit,backgrounds,decorations.pathreplacing}
\usepackage{algorithm}
\usepackage{algpseudocode}
\usepackage{listings}
\usepackage{xcolor}
\usepackage{tcolorbox}
\tcbuselibrary{skins,breakable,theorems}
\usepackage{fancyhdr}
\usepackage{titlesec}

% hyperref 放在最后加载
\usepackage[colorlinks=true,linkcolor=blue!70!black,citecolor=green!50!black,urlcolor=blue!60!black,bookmarks=true,bookmarksnumbered=true]{hyperref}

% ============ 颜色定义 ============
\definecolor{ugcolor}{RGB}{0,102,204}       % 本科生 - 蓝色
\definecolor{gradcolor}{RGB}{204,51,0}      % 研究生 - 红色
\definecolor{selfcolor}{RGB}{0,153,76}      % 自学拓展 - 绿色
\definecolor{codebg}{RGB}{245,245,245}
\definecolor{codeframe}{RGB}{200,200,200}
\definecolor{partcolor}{RGB}{0,51,102}
\definecolor{caseAcolor}{RGB}{0,100,180}    % 案例A - 深蓝
\definecolor{caseBcolor}{RGB}{180,60,0}     % 案例B - 深橙

% ============ 受众标记命令 ============
\newcommand{\ug}{\texorpdfstring{\textcolor{ugcolor}{\textbf{[本]}}}{[本]}}
\newcommand{\grad}{\texorpdfstring{\textcolor{gradcolor}{\textbf{[研]}}}{[研]}}
\newcommand{\self}{\texorpdfstring{\textcolor{selfcolor}{\textbf{[拓]}}}{[拓]}}

% ============ 案例标记命令 ============
\newcommand{\caseA}[1]{\textcolor{caseAcolor}{\textbf{【案例A-#1】}}}
\newcommand{\caseB}[1]{\textcolor{caseBcolor}{\textbf{【案例B-#1】}}}

% ============ 定理环境 ============
\theoremstyle{definition}
\newtheorem{definition}{定义}[chapter]
\newtheorem{example}{例}[chapter]
\newtheorem{remark}{注}[chapter]
\theoremstyle{plain}
\newtheorem{theorem}{定理}[chapter]
\newtheorem{proposition}{命题}[chapter]
\newtheorem{lemma}{引理}[chapter]

% ============ tcolorbox 环境 ============
\newtcolorbox{keypoint}[1][]{
  colback=blue!5!white,
  colframe=blue!60!black,
  title={\textbf{要点}},
  fonttitle=\bfseries,
  breakable,
  #1
}

\newtcolorbox{algbox}[1][]{
  colback=gray!5!white,
  colframe=gray!60!black,
  title={\textbf{算法概要}},
  fonttitle=\bfseries,
  breakable,
  #1
}

\newtcolorbox{appbox}[1][]{
  colback=green!5!white,
  colframe=green!60!black,
  title={\textbf{应用案例}},
  fonttitle=\bfseries,
  breakable,
  #1
}

\newtcolorbox{levelbox}[1][]{
  colback=orange!5!white,
  colframe=orange!60!black,
  title={\textbf{学习建议}},
  fonttitle=\bfseries,
  breakable,
  #1
}

% 案例A环境（多星观测）
\newtcolorbox{caseboxA}[1][]{
  colback=blue!3!white,
  colframe=caseAcolor,
  title={\textbf{案例A：多星协同对地观测}},
  fonttitle=\bfseries,
  breakable,
  #1
}

% 案例B环境（军事战术）
\newtcolorbox{caseboxB}[1][]{
  colback=orange!3!white,
  colframe=caseBcolor,
  title={\textbf{案例B：军事战术任务规划}},
  fonttitle=\bfseries,
  breakable,
  #1
}

% 讨论题环境
\newtcolorbox{discussbox}[1][]{
  colback=yellow!5!white,
  colframe=yellow!60!black,
  title={\textbf{讨论与思考}},
  fonttitle=\bfseries,
  breakable,
  #1
}

% ============ 代码环境 ============
\lstset{
  basicstyle=\ttfamily\small,
  backgroundcolor=\color{codebg},
  frame=single,
  rulecolor=\color{codeframe},
  breaklines=true,
  numbers=left,
  numberstyle=\tiny\color{gray},
  keywordstyle=\color{blue!70!black}\bfseries,
  commentstyle=\color{green!50!black},
  stringstyle=\color{red!60!black},
  showstringspaces=false
}

% ============ 页眉页脚 ============
\pagestyle{fancy}
\fancyhf{}
\fancyhead[LE]{\leftmark}
\fancyhead[RO]{\rightmark}
\fancyfoot[C]{\thepage}
\renewcommand{\headrulewidth}{0.4pt}

% ============ 章节标题格式 ============
\ctexset{
  chapter/format = \LARGE\bfseries\color{blue!70!black},
  section/format = \Large\bfseries,
  subsection/format = \large\bfseries
}

% ============ 算法环境中文化 ============
\floatname{algorithm}{算法}
\renewcommand{\algorithmicrequire}{\textbf{输入:}}
\renewcommand{\algorithmicensure}{\textbf{输出:}}

% ============ 文档信息 ============
\title{{\Huge\bfseries 任务规划算法体系}\\[0.8em]
{\Large 案例式教学版}\\[0.5em]
{\large 从经典方法到智能规划}\\[2em]
{\normalsize\color{gray} 本科生教学参考教材}}
\author{}
\date{\today}

\begin{document}

% ============ 封面 ============
\maketitle

% ============ 教学说明 ============
\thispagestyle{empty}
\vspace*{2cm}
\begin{center}
{\Large\bfseries 案例式教学说明}
\end{center}
\vspace{1cm}

\noindent 本书采用\textbf{案例式教学}方法，以两个核心案例贯穿全书：

\vspace{1em}

\begin{center}
\begin{tabular}{cp{10cm}}
\toprule
\textbf{案例线} & \textbf{内容描述} \\
\midrule
\textcolor{caseAcolor}{\textbf{案例A}} & \textbf{多星协同对地观测任务规划}：从单星确定性观测开始，逐步引入约束、不确定性、多星协同、对抗博弈等要素。 \\[0.8em]
\textcolor{caseBcolor}{\textbf{案例B}} & \textbf{军事战术任务规划}：从单平台侦察任务开始，逐步演化为多阶段任务分解、不完全信息决策、多平台协同、欺骗与隐蔽规划。 \\
\bottomrule
\end{tabular}
\end{center}

\vspace{1.5em}

\noindent\textbf{案例演进逻辑：}
\begin{itemize}[leftmargin=2em]
  \item \textbf{v1$\rightarrow$v2}：从"能做什么"到"在约束下做什么"——引入HTN分解与约束求解
  \item \textbf{v2$\rightarrow$v3}：从"确定世界"到"不确定世界"——引入概率规划与POMDP
  \item \textbf{v3$\rightarrow$v4}：从"单体决策"到"群体协同"——引入分布式规划与MARL
  \item \textbf{v4$\rightarrow$v5}：从"合作环境"到"对抗环境"——引入博弈规划与对抗规划
\end{itemize}

\vspace{1.5em}

\noindent\textbf{建议学习方式：}
\begin{enumerate}[leftmargin=2em]
  \item \textbf{课前}：阅读案例描述，思考预习问题（约30分钟）
  \item \textbf{课中}：参与案例讨论，学习算法原理，分析案例求解过程
  \item \textbf{课后}：完成编程练习，尝试改进算法，撰写案例分析报告
\end{enumerate}

\clearpage

% ============ 目录 ============
\tableofcontents

% ============ 前言 ============
\chapter*{前言}
\addcontentsline{toc}{chapter}{前言}

任务规划（Task Planning）是人工智能领域的核心研究方向之一，其目标是让智能体能够自主地制定行动方案以达成预设目标。本书采用\textbf{案例式教学}方法，以两个核心案例——"多星协同对地观测"和"军事战术任务规划"——贯穿全书七章内容，帮助学生在解决实际问题的过程中掌握任务规划的核心算法与方法。

\textbf{案例式教学的理念}：
\begin{itemize}[leftmargin=2em]
  \item 以真实问题为驱动，激发学习兴趣
  \item 通过案例演进揭示方法演化的内在逻辑
  \item 在讨论中深化理解，在实践中巩固知识
\end{itemize}

本书内容组织如下：
\begin{itemize}[leftmargin=2em]
  \item \textbf{第1章}：绪论与背景知识——介绍任务规划基本概念，回顾搜索方法，引入核心案例
  \item \textbf{第2章}：任务建模与经典规划——STRIPS/PDDL建模，前向搜索
  \item \textbf{第3章}：层次任务网络与约束规划——HTN分解，CSP/SAT/SMT约束求解
  \item \textbf{第4章}：不确定性与概率规划——MDP/POMDP，强化学习基础
  \item \textbf{第5章}：分布式与多智能体规划——MA-STRIPS，MAPF，MARL
  \item \textbf{第6章}：博弈规划与对抗规划——Stackelberg博弈，欺骗与隐蔽规划
  \item \textbf{第7章}：前沿方法与系统集成——LLM规划，神经符号方法，系统设计
\end{itemize}

\vspace{2em}
\hfill 编者

\hfill \today

% =============================================
% 第1章：绪论与背景知识
% =============================================
% 第1章：绪论与背景知识
\chapter{绪论与背景知识}

\begin{levelbox}
本章为全书的基础，介绍任务规划的基本概念、应用领域、搜索方法回顾，并引入贯穿全书的两个核心案例。建议学时：4学时。
\end{levelbox}

% =============================================
\section{任务规划的基本概念与应用领域}
% =============================================

\subsection{什么是任务规划}

\begin{definition}[任务规划]
\textbf{任务规划}（Task Planning）是指智能体根据当前状态、目标状态和可用动作，自动生成一个能够达成目标的动作序列或策略的过程。
\end{definition}

任务规划是人工智能的核心问题之一。早在1959年，Newell和Simon就提出了通用问题求解器（GPS），开创了自动规划研究的先河。经过数十年发展，任务规划已形成完整的理论体系和丰富的应用实践。

\subsection{任务规划的核心要素}

一个任务规划问题通常包含以下要素：

\begin{itemize}[leftmargin=2em]
  \item \textbf{状态空间} $\mathcal{S}$：所有可能状态的集合
  \item \textbf{初始状态} $s_0 \in \mathcal{S}$：问题的起始状态
  \item \textbf{目标} $G$：需要达成的条件，可以是目标状态 $s_g$ 或目标条件集合
  \item \textbf{动作集合} $\mathcal{A}$：智能体可执行的所有动作
  \item \textbf{状态转移函数} $\delta: \mathcal{S} \times \mathcal{A} \rightarrow \mathcal{S}$：描述动作对状态的影响
  \item \textbf{规划} $\pi$：一个动作序列 $\langle a_1, a_2, \ldots, a_n \rangle$，使得从 $s_0$ 出发依次执行后可达成目标
\end{itemize}

\subsection{任务规划的主要应用领域}

\begin{table}[htbp]
\centering
\caption{任务规划的典型应用领域}
\begin{tabular}{p{3cm}p{4cm}p{5cm}}
\toprule
\textbf{领域} & \textbf{典型问题} & \textbf{特点} \\
\midrule
航天航空 & 卫星观测调度、航迹规划 & 时间窗口约束、资源有限 \\
军事国防 & 作战任务规划、无人机协同 & 对抗性、不确定性高 \\
物流运输 & 车辆路径规划、仓储调度 & 大规模、实时性要求 \\
机器人 & 运动规划、任务分配 & 连续空间、多约束 \\
智能制造 & 生产调度、工艺规划 & 资源冲突、时序约束 \\
\bottomrule
\end{tabular}
\end{table}

% =============================================
\section{搜索方法回顾}
% =============================================

\begin{levelbox}
本节回顾搜索算法的基础知识，作为后续章节的预备知识。已熟悉搜索算法的读者可快速浏览。
\end{levelbox}

\subsection{状态空间搜索框架}

任务规划可以抽象为状态空间中的搜索问题。从初始状态出发，通过执行动作在状态空间中探索，直到找到满足目标的状态。

\begin{algbox}[通用状态空间搜索框架]
\begin{enumerate}
  \item 初始化：将初始状态 $s_0$ 放入开放列表（Open List）
  \item 循环：
  \begin{enumerate}
    \item 如果开放列表为空，返回失败
    \item 从开放列表中选择一个状态 $s$（选择策略决定搜索方式）
    \item 如果 $s$ 满足目标条件，返回解路径
    \item 将 $s$ 移入关闭列表（Closed List）
    \item 扩展 $s$：对每个可执行动作 $a$，计算后继状态 $s' = \delta(s, a)$
    \item 将新的后继状态加入开放列表
  \end{enumerate}
\end{enumerate}
\end{algbox}

\subsection{盲目搜索策略}

\textbf{广度优先搜索（BFS）}：按层次扩展，保证找到最短路径（动作数最少），但空间复杂度高。

\textbf{深度优先搜索（DFS）}：沿一条路径深入探索，空间复杂度低，但不保证最优性。

\textbf{迭代加深搜索（IDS）}：结合BFS的完备性和DFS的低空间复杂度。

\subsection{启发式搜索与A*算法}

\begin{keypoint}[A*算法核心思想]
A*算法使用评价函数 $f(n) = g(n) + h(n)$ 指导搜索：
\begin{itemize}
  \item $g(n)$：从初始状态到节点 $n$ 的实际代价
  \item $h(n)$：从节点 $n$ 到目标的估计代价（启发式函数）
  \item 当 $h(n)$ 满足\textbf{可纳性}（不高估实际代价）时，A*保证找到最优解
\end{itemize}
\end{keypoint}

\begin{definition}[可纳启发式]
启发式函数 $h$ 是\textbf{可纳的}（admissible），当且仅当对所有状态 $n$：
\[ h(n) \leq h^*(n) \]
其中 $h^*(n)$ 是从 $n$ 到目标的真实最优代价。
\end{definition}

\subsection{启发式函数设计}

好的启发式函数应该：
\begin{enumerate}[leftmargin=2em]
  \item 提供有效的搜索引导（接近真实代价）
  \item 计算代价低
  \item 满足可纳性（保证最优性）
\end{enumerate}

常见的启发式设计方法：
\begin{itemize}[leftmargin=2em]
  \item \textbf{松弛问题}：忽略部分约束，求解简化问题的最优解作为启发值
  \item \textbf{模式数据库}：预计算子问题的最优解，存储为查找表
  \item \textbf{地标启发式}：识别必须达成的子目标（地标），统计未完成数量
\end{itemize}

% =============================================
\section{核心案例介绍}
% =============================================

本书采用两个核心案例贯穿全部章节，通过案例的逐步演化展示不同规划方法的适用场景和演化逻辑。

\subsection{案例A：多星协同对地观测任务规划}

\begin{caseboxA}[案例A：多星协同对地观测任务规划]
\textbf{背景}：某国拥有由多颗遥感卫星组成的对地观测系统，需要对地面多个目标区域进行成像观测。

\textbf{问题}：如何调度各卫星的观测任务，在满足各种约束的条件下，最大化观测收益？

\textbf{案例版本演进}：
\begin{itemize}
  \item \textbf{v1}（第2章）：单星、确定性、无复杂约束 → PDDL建模
  \item \textbf{v2}（第3章）：增加时间窗口、能量、存储约束 → HTN + CSP
  \item \textbf{v3}（第4章）：引入云层遮挡不确定性 → MDP/POMDP
  \item \textbf{v4}（第5章）：多星协同覆盖、资源竞争 → 分布式规划
  \item \textbf{v5}（第6章）：敌方电子干扰、轨道机动对抗 → 博弈规划
\end{itemize}
\end{caseboxA}

\textbf{案例A数据}（简化版）：
\begin{itemize}[leftmargin=2em]
  \item 卫星数量：1颗（v1-v3），3颗（v4-v5）
  \item 目标区域：10个待观测点，各有优先级和观测收益
  \item 轨道周期：约90分钟，每轨可观测3-5个目标
  \item 主要约束：观测窗口、存储容量、能量、下传能力
\end{itemize}

\subsection{案例B：军事战术任务规划}

\begin{caseboxB}[案例B：军事战术任务规划]
\textbf{背景}：某战术分队需要执行侦察-打击-评估任务链，涉及无人机、地面平台等多种作战单元。

\textbf{问题}：如何规划各平台的任务序列，在不确定的战场环境中完成作战目标？

\textbf{案例版本演进}：
\begin{itemize}
  \item \textbf{v1}（第2章）：单无人机侦察任务 → 基本建模
  \item \textbf{v2}（第3章）：侦察-打击-评估任务链分解 → HTN分解
  \item \textbf{v3}（第4章）：战场迷雾、敌情不明 → POMDP
  \item \textbf{v4}（第5章）：有人-无人协同编队 → 多智能体规划
  \item \textbf{v5}（第6章）：欺骗机动、意图隐蔽 → 欺骗/隐蔽规划
\end{itemize}
\end{caseboxB}

\textbf{案例B数据}（简化版）：
\begin{itemize}[leftmargin=2em]
  \item 作战单元：1架无人机（v1-v3），1架有人机+2架无人机（v4-v5）
  \item 目标：3个高价值目标，需要侦察确认后实施打击
  \item 威胁：防空火力点、电子干扰
  \item 主要约束：航程、载弹量、通信链路、生存概率
\end{itemize}

\subsection{案例演进的逻辑说明}

\begin{table}[htbp]
\centering
\caption{案例版本演进的逻辑}
\begin{tabular}{ccp{8cm}}
\toprule
\textbf{演进} & \textbf{引入要素} & \textbf{为什么需要新方法？} \\
\midrule
v1→v2 & 复杂约束 & 前一方法难以处理时间、资源等硬约束 \\
v2→v3 & 不确定性 & 现实中观测/感知存在不确定性，需要概率建模 \\
v3→v4 & 多智能体 & 单体能力有限，需要多个智能体协同完成任务 \\
v4→v5 & 对抗博弈 & 存在敌方对手，需要考虑博弈和对抗策略 \\
\bottomrule
\end{tabular}
\end{table}

% =============================================
\section{本书结构与学习方法}
% =============================================

\subsection{方法论视角：数据-模型-算法的融合}

当代任务规划研究呈现出\textbf{数据驱动}与\textbf{模型驱动}深度融合的趋势，\textbf{算法}作为桥梁实现二者的协同优化。这一融合视角贯穿本书各章：

\begin{itemize}[leftmargin=2em]
  \item \textbf{模型驱动}（第2-3章）：PDDL/HTN等形式化模型，利用领域知识结构化表达
  \item \textbf{数据驱动}（第4-5章）：强化学习、模仿学习等从交互数据中学习策略
  \item \textbf{深度融合}（第7章）：LLM+规划器、神经符号方法、世界模型等融合范式
\end{itemize}

华中科技大学王红卫、祁超团队在不确定HTN规划领域的研究，为模型与算法的融合提供了典型范例——将领域任务分解知识（模型）嵌入规划搜索与协调机制（算法），有效应对应急响应等复杂管理任务。

\subsection{全书章节结构}

\begin{table}[htbp]
\centering
\caption{全书章节与案例对应关系}
\begin{tabular}{clcc}
\toprule
\textbf{章} & \textbf{内容} & \textbf{案例A} & \textbf{案例B} \\
\midrule
1 & 绪论与背景知识 & 介绍 & 介绍 \\
2 & 任务建模与经典规划 & v1 & v1 \\
3 & 层次任务网络与约束规划 & v2 & v2 \\
4 & 不确定性与概率规划 & v3 & v3 \\
5 & 分布式与多智能体规划 & v4 & v4 \\
6 & 博弈规划与对抗规划 & v5 & v5 \\
7 & 前沿方法与系统集成 & 综合 & 综合 \\
\bottomrule
\end{tabular}
\end{table}

\subsection{建议学习方式}

\begin{enumerate}[leftmargin=2em]
  \item \textbf{课前预习}：阅读案例描述，思考如何解决案例中的问题
  \item \textbf{课堂学习}：
  \begin{itemize}
    \item 讨论案例问题的抽象和建模
    \item 学习相关算法的原理
    \item 分析算法在案例上的应用
  \end{itemize}
  \item \textbf{课后巩固}：
  \begin{itemize}
    \item 完成编程练习，实现核心算法
    \item 在案例数据上运行和调试
    \item 撰写案例分析报告
  \end{itemize}
\end{enumerate}

% =============================================
\section{本章小结与讨论}
% =============================================

\subsection*{本章小结}

\begin{itemize}[leftmargin=2em]
  \item 任务规划是让智能体自主生成行动方案的核心AI问题
  \item 状态空间搜索是任务规划的基础框架
  \item A*算法通过启发式函数有效引导搜索
  \item 本书以"多星观测"和"军事战术"两个案例贯穿全书
\end{itemize}

\begin{discussbox}
\textbf{讨论题}：
\begin{enumerate}
  \item 除了本章提到的应用领域，你还能想到哪些任务规划的应用场景？
  \item 对于案例A（卫星观测），你认为最大的挑战是什么？如果只用A*搜索能解决吗？
  \item 案例B涉及军事应用，你认为任务规划在军事领域应用需要注意哪些伦理问题？
\end{enumerate}
\end{discussbox}


% =============================================
% 第2章：任务建模与经典规划
% =============================================
% 第2章：任务建模与经典规划
\chapter{任务建模与经典规划}

\begin{levelbox}
本章介绍任务规划的形式化建模方法（STRIPS/PDDL）和基于前向搜索的经典规划算法。通过案例A-v1和案例B-v1学习如何将实际问题抽象为规划问题。建议学时：6学时。
\end{levelbox}

% =============================================
\section{\caseA{v1}单星观测任务：问题抽象}
% =============================================

\begin{caseboxA}[案例A-v1：单星确定性观测任务]
\textbf{场景设定}：
\begin{itemize}
  \item 1颗遥感卫星，需观测地面5个目标点
  \item 每个目标有观测收益值
  \item 卫星按固定轨道运行，每轨可观测若干目标
  \item 目标：选择观测序列，最大化总收益
\end{itemize}

\textbf{简化假设}（v1版本）：
\begin{itemize}
  \item 忽略时间窗口约束（假设所有目标均可观测）
  \item 忽略存储和能量约束
  \item 确定性环境（观测必定成功）
\end{itemize}

\textbf{预习思考}：
\begin{enumerate}
  \item 如何定义卫星观测问题的状态？
  \item 需要哪些动作？动作的前置条件和效果是什么？
  \item 如何表示目标条件？
\end{enumerate}
\end{caseboxA}

\subsection{问题抽象}

将卫星观测问题抽象为规划问题：

\textbf{状态表示}：
\begin{itemize}[leftmargin=2em]
  \item 每个目标点的观测状态：已观测/未观测
  \item 卫星的当前位置（可离散化为轨道位置编号）
\end{itemize}

\textbf{动作定义}：
\begin{itemize}[leftmargin=2em]
  \item \texttt{observe(target)}：观测目标target
  \item \texttt{move(pos1, pos2)}：卫星从位置pos1移动到pos2
\end{itemize}

\textbf{目标}：所有高优先级目标均被观测

% =============================================
\section{STRIPS/PDDL建模语言}
% =============================================

\subsection{STRIPS表示}

STRIPS（Stanford Research Institute Problem Solver）是最经典的规划问题表示语言。

\begin{definition}[STRIPS规划问题]
一个STRIPS规划问题由四元组 $\langle P, O, I, G \rangle$ 定义：
\begin{itemize}
  \item $P$：命题（谓词）集合
  \item $O$：操作（动作）集合，每个操作 $o$ 包含：
    \begin{itemize}
      \item $\text{Precond}(o)$：前置条件（命题集合）
      \item $\text{Add}(o)$：添加效果（动作后为真的命题）
      \item $\text{Del}(o)$：删除效果（动作后为假的命题）
    \end{itemize}
  \item $I \subseteq P$：初始状态
  \item $G \subseteq P$：目标条件
\end{itemize}
\end{definition}

\subsection{PDDL建模}

PDDL（Planning Domain Definition Language）是STRIPS的标准化扩展，是国际规划竞赛（IPC）的标准语言。

\begin{keypoint}[PDDL的两个文件]
\begin{itemize}
  \item \textbf{领域文件}（Domain File）：定义谓词和动作模式，与具体问题实例无关
  \item \textbf{问题文件}（Problem File）：定义具体的对象、初始状态和目标
\end{itemize}
\end{keypoint}

\subsection{案例A-v1的PDDL建模}

\textbf{领域文件}（satellite-domain.pddl）：

\begin{lstlisting}[language=Lisp,caption=卫星观测领域定义]
(define (domain satellite-observe)
  (:requirements :strips)

  (:predicates
    (satellite ?s)           ; s是卫星
    (target ?t)              ; t是目标
    (position ?p)            ; p是位置
    (at ?s ?p)               ; 卫星s在位置p
    (visible ?t ?p)          ; 目标t在位置p可见
    (observed ?t)            ; 目标t已被观测
  )

  (:action observe
    :parameters (?s ?t ?p)
    :precondition (and (satellite ?s) (target ?t) (position ?p)
                       (at ?s ?p) (visible ?t ?p) (not (observed ?t)))
    :effect (observed ?t)
  )

  (:action move
    :parameters (?s ?from ?to)
    :precondition (and (satellite ?s) (position ?from) (position ?to)
                       (at ?s ?from))
    :effect (and (at ?s ?to) (not (at ?s ?from)))
  )
)
\end{lstlisting}

\textbf{问题文件}（satellite-problem.pddl）：

\begin{lstlisting}[language=Lisp,caption=卫星观测问题实例]
(define (problem observe-5targets)
  (:domain satellite-observe)

  (:objects
    sat1 - satellite
    t1 t2 t3 t4 t5 - target
    p1 p2 p3 p4 p5 - position
  )

  (:init
    (satellite sat1)
    (target t1) (target t2) (target t3) (target t4) (target t5)
    (position p1) (position p2) (position p3) (position p4) (position p5)
    (at sat1 p1)
    (visible t1 p1) (visible t2 p2) (visible t3 p3)
    (visible t4 p4) (visible t5 p5)
  )

  (:goal (and (observed t1) (observed t2) (observed t3)))
)
\end{lstlisting}

% =============================================
\section{前向搜索与启发式规划}
% =============================================

\subsection{前向状态空间搜索}

前向搜索从初始状态出发，逐步应用动作，直到达到目标状态。

\begin{algorithm}[htbp]
\caption{前向状态空间搜索}
\begin{algorithmic}[1]
\Require 规划问题 $\langle P, O, I, G \rangle$
\Ensure 动作序列或失败
\State $\text{Open} \gets \{I\}$，$\text{Closed} \gets \emptyset$
\While{$\text{Open} \neq \emptyset$}
  \State $s \gets$ 从Open中选择状态（按启发式值）
  \If{$G \subseteq s$}
    \State \Return 从$I$到$s$的动作序列
  \EndIf
  \State $\text{Closed} \gets \text{Closed} \cup \{s\}$
  \For{每个可应用的动作 $a \in O$}
    \State $s' \gets \text{apply}(s, a)$
    \If{$s' \notin \text{Closed}$}
      \State 将 $s'$ 加入 Open
    \EndIf
  \EndFor
\EndWhile
\State \Return 失败
\end{algorithmic}
\end{algorithm}

\subsection{规划启发式}

\subsubsection{删除松弛启发式}

\textbf{删除松弛}（Delete Relaxation）：忽略动作的删除效果，只保留添加效果。松弛问题更容易求解，其最优解代价是原问题的下界。

\begin{itemize}[leftmargin=2em]
  \item $h^{\text{add}}$：松弛问题中所有子目标代价之和
  \item $h^{\text{max}}$：松弛问题中代价最大的子目标
  \item $h^{\text{FF}}$：基于规划图的松弛启发式，FF规划器使用
\end{itemize}

\subsubsection{Fast Downward与LAMA}

\textbf{Fast Downward}是现代最成功的规划器之一，采用：
\begin{itemize}[leftmargin=2em]
  \item SAS+表示：多值变量替代命题
  \item 因果图分析：识别变量间的依赖关系
  \item 多种启发式：$h^{\text{FF}}$、因果图启发式、地标启发式
\end{itemize}

\textbf{LAMA}基于地标启发式，多次获得IPC冠军。

% =============================================
\section{\caseB{v1}单平台侦察任务建模}
% =============================================

\begin{caseboxB}[案例B-v1：单无人机侦察任务]
\textbf{场景设定}：
\begin{itemize}
  \item 1架无人机，需侦察3个目标区域
  \item 无人机从基地起飞，完成侦察后返回
  \item 每个目标区域有侦察价值
  \item 目标：规划飞行路线，完成所有侦察任务
\end{itemize}

\textbf{简化假设}（v1版本）：
\begin{itemize}
  \item 忽略航程约束
  \item 忽略敌方威胁
  \item 确定性环境
\end{itemize}
\end{caseboxB}

\subsection{问题建模}

\textbf{状态表示}：
\begin{itemize}[leftmargin=2em]
  \item 无人机位置
  \item 各目标的侦察状态（已侦察/未侦察）
\end{itemize}

\textbf{动作}：
\begin{itemize}[leftmargin=2em]
  \item \texttt{fly(from, to)}：飞行
  \item \texttt{recon(target)}：侦察目标
  \item \texttt{return\_base()}：返回基地
\end{itemize}

\subsection{PDDL建模}

\begin{lstlisting}[language=Lisp,caption=无人机侦察领域定义]
(define (domain uav-recon)
  (:requirements :strips)

  (:predicates
    (uav ?u)
    (location ?l)
    (target ?t)
    (at ?u ?l)
    (target-at ?t ?l)
    (reconned ?t)
    (at-base ?u)
  )

  (:action fly
    :parameters (?u ?from ?to)
    :precondition (and (uav ?u) (location ?from) (location ?to)
                       (at ?u ?from))
    :effect (and (at ?u ?to) (not (at ?u ?from)) (not (at-base ?u)))
  )

  (:action recon
    :parameters (?u ?t ?l)
    :precondition (and (uav ?u) (target ?t) (location ?l)
                       (at ?u ?l) (target-at ?t ?l) (not (reconned ?t)))
    :effect (reconned ?t)
  )

  (:action return-base
    :parameters (?u ?l)
    :precondition (and (uav ?u) (location ?l) (at ?u ?l))
    :effect (at-base ?u)
  )
)
\end{lstlisting}

% =============================================
\section{本章小结与讨论}
% =============================================

\subsection*{本章小结}

\begin{itemize}[leftmargin=2em]
  \item STRIPS/PDDL是任务规划的标准建模语言
  \item 前向状态空间搜索配合启发式函数是经典规划的主流方法
  \item 删除松弛是设计可纳启发式的有效技术
  \item Fast Downward和LAMA是现代高效规划器的代表
\end{itemize}

\begin{discussbox}
\textbf{讨论题}：
\begin{enumerate}
  \item 案例A-v1的PDDL建模忽略了哪些实际约束？这些约束能用基本PDDL表达吗？
  \item 如果目标数量从5个增加到100个，搜索效率会如何变化？有什么改进方法？
  \item 案例B的"返回基地"是否必须作为显式目标？如果不返回会怎样？
  \item 尝试为案例A或案例B设计一个简单的启发式函数，分析其可纳性。
\end{enumerate}

\textbf{编程练习}：
\begin{enumerate}
  \item 使用Fast Downward求解案例A-v1的PDDL问题
  \item 增加目标数量，测试求解时间的变化趋势
  \item 修改目标条件（部分观测即可），比较规划结果的差异
\end{enumerate}
\end{discussbox}


% =============================================
% 第3章：层次任务网络与约束规划
% =============================================
% 第3章：层次任务网络与约束规划
\chapter{层次任务网络与约束规划}

\begin{levelbox}
本章介绍层次任务网络（HTN）规划和基于约束的规划方法。通过案例v2版本学习如何处理任务分解和复杂约束。建议学时：8学时。
\end{levelbox}

% =============================================
\section{\caseB{v2}多阶段战术任务分解}
% =============================================

\begin{caseboxB}[案例B-v2：侦察-打击-评估任务链]
\textbf{场景演进}：v1只考虑单一侦察任务，v2引入完整的任务链——"侦察-打击-评估"（ISR-Strike-BDA）。

\textbf{新增要素}：
\begin{itemize}
  \item 任务必须按照"先侦察、后打击、再评估"的顺序执行
  \item 打击需要基于侦察结果选择武器和攻击方式
  \item 打击后需要进行战斗损伤评估（BDA）
\end{itemize}

\textbf{为什么v1方法不够？}
\begin{itemize}
  \item 经典PDDL难以自然表达任务的层次结构
  \item 复杂任务序列的建模变得困难
  \item 领域专家知识（如任务分解规则）难以融入
\end{itemize}

\textbf{预习思考}：
\begin{enumerate}
  \item "侦察-打击-评估"可以如何分解为子任务？
  \item 子任务之间有什么依赖关系？
  \item 如何利用领域知识指导规划？
\end{enumerate}
\end{caseboxB}

% =============================================
\section{HTN规划：SHOP/SHOP2/HDDL}
% =============================================

\subsection{HTN基本概念}

\begin{definition}[层次任务网络（HTN）]
HTN规划采用自顶向下的任务分解策略：
\begin{itemize}
  \item \textbf{原子任务}（Primitive Tasks）：可直接执行的动作
  \item \textbf{复合任务}（Compound Tasks）：需要分解的高层任务
  \item \textbf{方法}（Methods）：定义复合任务如何分解为子任务
\end{itemize}
\end{definition}

\begin{keypoint}[HTN vs 经典规划]
\begin{itemize}
  \item 经典规划：搜索达成目标的动作序列
  \item HTN规划：搜索将初始任务分解为原子任务的方法序列
  \item HTN天然支持领域知识的编码（通过方法定义）
\end{itemize}
\end{keypoint}

\subsection{SHOP/SHOP2系统}

SHOP（Simple Hierarchical Ordered Planner）是最著名的HTN规划器：

\begin{itemize}[leftmargin=2em]
  \item 使用\textbf{全序分解}：子任务有严格的线性顺序
  \item 支持在前置条件中使用任意Lisp表达式
  \item 在实际应用中表现优异（军事规划、Web服务组合等）
\end{itemize}

\subsection{案例B-v2的HTN建模}

\begin{lstlisting}[caption=案例B-v2的HTN方法定义（SHOP风格伪代码）]
; 复合任务：执行完整的打击任务链
(:method (attack-mission ?target)
  ; 方法1：标准流程
  :precondition (enemy-target ?target)
  :subtasks (
    (recon ?target)           ; 先侦察
    (strike ?target)          ; 后打击
    (damage-assess ?target)   ; 再评估
  )
)

; 复合任务：侦察
(:method (recon ?target)
  ; 方法1：使用无人机侦察
  :precondition (uav-available)
  :subtasks (
    (fly-to ?target)
    (take-images ?target)
    (transmit-data ?target)
  )
)

; 复合任务：打击
(:method (strike ?target)
  ; 方法1：精确制导打击（侦察已确认高价值目标）
  :precondition (and (reconned ?target) (high-value ?target))
  :subtasks (
    (select-weapon precision-guided)
    (fly-to-attack-position ?target)
    (launch-weapon ?target)
  )
  ; 方法2：区域压制（目标位置不精确）
  :precondition (and (reconned ?target) (not (precise-location ?target)))
  :subtasks (
    (select-weapon area-suppression)
    (fly-to-attack-position ?target)
    (launch-weapon ?target)
  )
)
\end{lstlisting}

\subsection{不确定HTN规划}

实际应用中，规划环境往往具有不确定性。华中科技大学王红卫教授、祁超副教授团队系统研究了不确定层次任务网络规划问题。

\begin{keypoint}[不确定HTN的三类不确定性（王红卫等, 2016）]
\begin{enumerate}
  \item \textbf{状态不确定}：初始状态或中间状态信息不完全，需要基于信念状态规划
  \item \textbf{动作效果不确定}：动作执行结果具有随机性，需要考虑多种可能后果
  \item \textbf{任务分解不确定}：方法选择本身受到不确定因素影响
\end{enumerate}
该研究将HTN规划与概率推理、条件规划等技术相结合，为应急响应等复杂管理任务提供了理论基础。
\end{keypoint}

\subsection{HDDL表示语言}

HDDL（Hierarchical Domain Definition Language）是HTN的标准化语言，扩展了PDDL以支持层次任务。

\begin{lstlisting}[language=Lisp,caption=HDDL任务定义示例]
(:task attack-mission
  :parameters (?target - location)
)

(:method m-attack-standard
  :parameters (?target - location)
  :task (attack-mission ?target)
  :precondition (enemy-target ?target)
  :ordered-subtasks (and
    (recon ?target)
    (strike ?target)
    (damage-assess ?target)
  )
)
\end{lstlisting}

% =============================================
\section{\caseA{v2}观测窗口约束下的任务分解}
% =============================================

\begin{caseboxA}[案例A-v2：带约束的卫星观测调度]
\textbf{场景演进}：v1忽略了所有约束，v2引入现实约束。

\textbf{新增约束}：
\begin{itemize}
  \item \textbf{时间窗口}：每个目标只在特定时段内可见
  \item \textbf{能量约束}：成像和数传消耗电能，需考虑能量平衡
  \item \textbf{存储约束}：星上存储有限，需及时下传数据
  \item \textbf{姿态约束}：相邻观测需要足够的姿态机动时间
\end{itemize}

\textbf{为什么需要约束求解？}
\begin{itemize}
  \item 时间窗口约束使问题成为时序约束满足问题
  \item 资源约束（能量、存储）需要累积计算
  \item 纯搜索方法难以高效处理这些约束
\end{itemize}
\end{caseboxA}

\subsection{HTN分解策略}

卫星观测调度可以分解为层次结构：

\begin{enumerate}[leftmargin=2em]
  \item \textbf{顶层任务}：完成观测任务集
  \item \textbf{轨道级分解}：为每个轨道周期安排观测
  \item \textbf{目标级分解}：为每个目标安排具体观测动作
  \item \textbf{原子动作}：姿态机动、成像、数据下传
\end{enumerate}

% =============================================
\section{约束满足规划（CSP/SAT/SMT）}
% =============================================

\subsection{CSP规划}

\begin{definition}[约束满足问题（CSP）]
CSP由三元组 $\langle X, D, C \rangle$ 定义：
\begin{itemize}
  \item $X = \{x_1, \ldots, x_n\}$：变量集合
  \item $D = \{D_1, \ldots, D_n\}$：每个变量的值域
  \item $C$：约束集合，定义变量间的合法组合
\end{itemize}
\end{definition}

将规划问题编码为CSP：
\begin{itemize}[leftmargin=2em]
  \item 变量：每个时间步的状态和动作
  \item 约束：初始状态约束、目标约束、动作前置条件和效果约束
\end{itemize}

\subsection{SAT规划}

\textbf{SAT规划}将规划问题编码为命题逻辑公式：

\begin{algbox}[SAT规划框架]
\begin{enumerate}
  \item 设定时间步上限 $T$
  \item 编码：
    \begin{itemize}
      \item 状态变量：$\text{holds}(p, t)$ 表示命题 $p$ 在时刻 $t$ 为真
      \item 动作变量：$\text{do}(a, t)$ 表示在时刻 $t$ 执行动作 $a$
      \item 初始/目标约束、因果约束、帧公理、互斥约束
    \end{itemize}
  \item 调用SAT求解器
  \item 若可满足，提取规划；否则增大 $T$ 重试
\end{enumerate}
\end{algbox}

\subsection{SMT规划}

SMT（Satisfiability Modulo Theories）扩展了SAT，支持：
\begin{itemize}[leftmargin=2em]
  \item 整数/实数算术约束
  \item 数组和位向量理论
  \item 适合处理数值变量（时间、资源）
\end{itemize}

\begin{example}[案例A-v2的SMT编码]
\begin{align*}
& \text{变量：} t_i \in \mathbb{R} \text{（目标}i\text{的观测时间）} \\
& \text{时间窗口约束：} w_i^{\text{start}} \leq t_i \leq w_i^{\text{end}} \\
& \text{姿态机动约束：} |t_i - t_j| \geq \tau_{\text{slew}} \text{ 或 } i,j\text{不相邻观测} \\
& \text{能量约束：} \sum_{i \text{ 被选中}} e_i \leq E_{\text{max}}
\end{align*}
\end{example}

% =============================================
\section{时态规划与资源约束}
% =============================================

\subsection{PDDL 2.1时态扩展}

PDDL 2.1引入\textbf{持续性动作}（Durative Actions）：

\begin{lstlisting}[language=Lisp,caption=持续性动作示例]
(:durative-action observe
  :parameters (?sat ?target)
  :duration (= ?duration 10)  ; 观测持续10个时间单位
  :condition (and
    (at start (visible ?target))
    (over all (has-power ?sat))
  )
  :effect (and
    (at end (observed ?target))
    (at start (decrease (energy ?sat) 5))
  )
)
\end{lstlisting}

\subsection{资源约束处理}

\begin{keypoint}[资源约束的处理方法]
\begin{enumerate}
  \item \textbf{PDDL数值扩展}：使用 \texttt{:numeric-fluents} 定义资源变量
  \item \textbf{约束传播}：在搜索过程中动态检查资源约束
  \item \textbf{LP松弛}：用线性规划检查资源可行性
\end{enumerate}
\end{keypoint}

% =============================================
\section{本章小结与讨论}
% =============================================

\subsection*{本章小结}

\begin{itemize}[leftmargin=2em]
  \item HTN通过任务分解捕获领域的层次结构
  \item SHOP/SHOP2是实际应用中最成功的HTN系统
  \item CSP/SAT/SMT将规划转化为约束求解问题
  \item 时态规划和资源约束扩展了经典规划的表达能力
\end{itemize}

\begin{discussbox}
\textbf{讨论题}：
\begin{enumerate}
  \item 案例B的"侦察-打击-评估"任务链是否必须严格顺序执行？有没有可以并行的部分？
  \item HTN方法需要领域专家定义，这是优点还是缺点？与机器学习自动获取知识相比呢？
  \item 案例A-v2中，如果时间窗口相互冲突（同时只能观测一个目标），应该如何处理？
  \item SAT规划和CSP规划各有什么优缺点？什么情况下选择哪种方法？
\end{enumerate}

\textbf{编程练习}：
\begin{enumerate}
  \item 使用PANDA或其他HTN规划器求解案例B-v2
  \item 将案例A-v2编码为SMT问题，使用Z3求解器求解
  \item 比较HTN规划和SAT规划在案例A-v2上的求解效率
\end{enumerate}
\end{discussbox}


% =============================================
% 第4章：不确定性与概率规划
% =============================================
% 第4章：不确定性与概率规划
\chapter{不确定性与概率规划}

\begin{levelbox}
本章介绍在不确定环境下的规划方法，包括条件规划、MDP/POMDP和强化学习基础。通过案例v3版本学习如何处理观测和状态不确定性。建议学时：8学时。
\end{levelbox}

% =============================================
\section{\caseA{v3}云层遮挡下的观测规划}
% =============================================

\begin{caseboxA}[案例A-v3：不确定性观测环境]
\textbf{场景演进}：v2假设观测必定成功，v3引入不确定性。

\textbf{新增不确定因素}：
\begin{itemize}
  \item \textbf{云层遮挡}：观测成功与否取决于云量，概率可预测但不确定
  \item \textbf{设备故障}：成像设备有一定故障概率
  \item \textbf{目标移动}：部分目标可能在观测窗口内移动
\end{itemize}

\textbf{为什么v2方法不够？}
\begin{itemize}
  \item 确定性规划假设动作效果完全可预测
  \item 无法处理"观测失败后重新规划"的场景
  \item 需要概率建模和期望收益优化
\end{itemize}

\textbf{预习思考}：
\begin{enumerate}
  \item 如何建模"观测可能失败"这一不确定性？
  \item 是否应该为每个目标安排备选观测窗口？
  \item 如何在不确定性下最大化期望观测收益？
\end{enumerate}
\end{caseboxA}

% =============================================
\section{条件规划与一致性规划}
% =============================================

\subsection{非确定性动作}

\begin{definition}[非确定性动作]
非确定性动作的效果不唯一，执行后可能到达多个可能状态之一：
\[ \text{Effects}(a) = \{e_1, e_2, \ldots, e_k\} \]
具体到达哪个状态在执行前未知。
\end{definition}

\subsection{条件规划}

\textbf{条件规划}（Contingent Planning）允许在执行过程中进行观测和分支。

\begin{keypoint}[条件规划的输出]
条件规划的输出不是线性动作序列，而是一棵\textbf{条件规划树}：
\begin{itemize}
  \item 内部节点：感知动作，根据观测结果分支
  \item 叶节点：目标状态或失败
  \item 执行时根据实际观测选择分支
\end{itemize}
\end{keypoint}

\begin{example}[案例A的条件规划]
\begin{verbatim}
1. 尝试观测目标T1
2. 感知：观测是否成功？
   - 如果成功：继续观测T2
   - 如果失败（云层遮挡）：
     - 等待下一观测窗口
     - 或转向备选目标T1'
\end{verbatim}
\end{example}

\subsection{一致性规划}

\textbf{一致性规划}（Conformant Planning）：在不确定初始状态下，寻找一个在所有可能初始状态下都能达成目标的动作序列。

\begin{itemize}[leftmargin=2em]
  \item 不依赖执行时观测
  \item 在信念空间（Belief Space）中搜索
  \item 计算复杂度高（EXPSPACE-完全）
\end{itemize}

\subsection{HTN规划中的不确定性}

值得注意的是，不确定性同样存在于层次任务网络规划中。华中科技大学王红卫教授、祁超教授团队将不确定HTN规划中的不确定性归纳为三类：\textbf{状态不确定}、\textbf{动作效果不确定}和\textbf{任务分解不确定}（参见Liu等人在Knowledge-Based Systems 2016发表的论文）。这一分类框架将第3章的HTN方法与本章的概率规划方法有机关联——HTN的层次分解结构可与概率推理相结合，为应急响应等不确定性环境下的复杂任务提供更具表达力的规划模型。

% =============================================
\section{MDP/POMDP框架}
% =============================================

\subsection{马尔可夫决策过程（MDP）}

\begin{definition}[MDP]
MDP由五元组 $\langle S, A, T, R, \gamma \rangle$ 定义：
\begin{itemize}
  \item $S$：状态空间
  \item $A$：动作空间
  \item $T(s'|s,a)$：状态转移概率
  \item $R(s,a)$：即时奖励函数
  \item $\gamma \in [0,1)$：折扣因子
\end{itemize}
\end{definition}

\textbf{目标}：找到策略 $\pi: S \rightarrow A$，最大化期望累积奖励：
\[ V^\pi(s) = \mathbb{E}\left[\sum_{t=0}^{\infty} \gamma^t R(s_t, \pi(s_t)) \mid s_0 = s\right] \]

\subsection{值迭代与策略迭代}

\begin{algbox}[值迭代算法]
\begin{enumerate}
  \item 初始化 $V(s) = 0$ 对所有 $s$
  \item 重复直到收敛：
  \[ V(s) \leftarrow \max_{a \in A} \left[ R(s,a) + \gamma \sum_{s'} T(s'|s,a) V(s') \right] \]
  \item 提取最优策略：
  \[ \pi^*(s) = \arg\max_{a} \left[ R(s,a) + \gamma \sum_{s'} T(s'|s,a) V(s') \right] \]
\end{enumerate}
\end{algbox}

\subsection{部分可观测MDP（POMDP）}

\begin{definition}[POMDP]
POMDP扩展MDP，增加观测模型：
\begin{itemize}
  \item $\Omega$：观测空间
  \item $O(o|s',a)$：在执行动作 $a$ 到达状态 $s'$ 后观测到 $o$ 的概率
\end{itemize}
\end{definition}

POMDP中智能体不能直接观测状态，而是维护一个\textbf{信念状态}（Belief State）$b$——关于当前状态的概率分布。

\subsection{案例A-v3的POMDP建模}

\begin{example}[云层遮挡的POMDP模型]
\textbf{状态}：$s = (\text{已观测目标集合}, \text{各目标云量状态})$

\textbf{动作}：$a = \text{observe}(t)$ 或 $\text{wait}$

\textbf{转移}：
\begin{itemize}
  \item 云量状态随时间演化（马尔可夫过程）
  \item 观测成功概率取决于云量
\end{itemize}

\textbf{观测}：
\begin{itemize}
  \item 执行observe后观测到"成功"或"失败"
  \item 云量状态不直接可观测，需通过气象预报推断
\end{itemize}

\textbf{奖励}：成功观测目标获得奖励，与目标优先级成正比
\end{example}

% =============================================
\section{\caseB{v3}战场迷雾下的决策}
% =============================================

\begin{caseboxB}[案例B-v3：不完全信息作战]
\textbf{场景演进}：v2假设敌情完全已知，v3引入"战场迷雾"。

\textbf{新增不确定性}：
\begin{itemize}
  \item \textbf{敌方位置不确定}：只知道敌方可能在某些区域
  \item \textbf{敌方意图不明}：不知道敌方是否会机动
  \item \textbf{侦察结果不完全}：侦察可能漏报或误报
\end{itemize}

\textbf{POMDP建模}：
\begin{itemize}
  \item 状态包含敌方真实位置和状态（不可直接观测）
  \item 侦察动作提供关于敌情的噪声观测
  \item 需要在信念空间中规划
\end{itemize}
\end{caseboxB}

% =============================================
\section{强化学习与世界模型}
% =============================================

\subsection{强化学习基础}

当转移概率 $T$ 和奖励函数 $R$ 未知时，可以通过\textbf{强化学习}（RL）从交互经验中学习策略。

\begin{itemize}[leftmargin=2em]
  \item \textbf{基于值的方法}：学习值函数 $Q(s,a)$，如Q-learning、DQN
  \item \textbf{基于策略的方法}：直接学习策略参数，如REINFORCE、PPO
  \item \textbf{Actor-Critic}：同时学习策略和值函数
\end{itemize}

\subsection{基于模型的强化学习}

\begin{keypoint}[基于模型的RL与规划]
基于模型的强化学习先学习环境模型（世界模型），再基于模型进行规划：
\begin{enumerate}
  \item 从经验中学习转移模型 $\hat{T}$ 和奖励模型 $\hat{R}$
  \item 在学习的模型中进行规划（如值迭代、蒙特卡洛树搜索）
  \item 执行规划结果，收集新经验
\end{enumerate}
\end{keypoint}

\subsection{MuZero与Dreamer}

\textbf{MuZero}（DeepMind, 2020）：
\begin{itemize}[leftmargin=2em]
  \item 学习隐状态表示和动态模型
  \item 在隐空间中进行蒙特卡洛树搜索
  \item 不需要环境的完整规则
  \item 在围棋、Atari等任务上达到超人水平
\end{itemize}

\textbf{Dreamer系列}（Hafner et al.）：
\begin{itemize}[leftmargin=2em]
  \item 学习世界模型（编码器+动态模型+解码器）
  \item 在"想象"（世界模型中的模拟）中训练策略
  \item DreamerV3在多种任务上表现优异
\end{itemize}

% =============================================
\section{本章小结与讨论}
% =============================================

\subsection*{本章小结}

\begin{itemize}[leftmargin=2em]
  \item 条件规划通过感知和分支处理不确定性
  \item MDP/POMDP是不确定性规划的数学框架
  \item 强化学习可以在未知环境中学习策略
  \item 世界模型（MuZero、Dreamer）结合了学习和规划
\end{itemize}

\begin{discussbox}
\textbf{讨论题}：
\begin{enumerate}
  \item 案例A-v3中，如果云量预报完全准确，问题是否退化为确定性规划？
  \item POMDP的精确求解是PSPACE-完全的，实际中如何处理大规模POMDP？
  \item 案例B-v3中，侦察误报（将非目标判为目标）和漏报（未发现目标）哪个更危险？如何在建模中体现？
  \item MuZero和传统规划器有什么本质区别？各有什么优缺点？
\end{enumerate}

\textbf{编程练习}：
\begin{enumerate}
  \item 将案例A-v3简化为小规模MDP，使用值迭代求解
  \item 使用POMDP-solve或SARSOP求解案例A-v3的POMDP
  \item 比较MDP策略和条件规划在模拟环境中的表现
\end{enumerate}
\end{discussbox}


% =============================================
% 第5章：分布式与多智能体规划
% =============================================
% 第5章：分布式与多智能体规划
\chapter{分布式与多智能体规划}

\begin{levelbox}
本章介绍多智能体环境下的分布式规划方法，包括MA-STRIPS、多智能体路径规划（MAPF）和多智能体强化学习（MARL）。通过案例v4版本学习如何处理多智能体协同。建议学时：8学时。
\end{levelbox}

% =============================================
\section{\caseA{v4}多星协同覆盖观测}
% =============================================

\begin{caseboxA}[案例A-v4：多星协同任务规划]
\textbf{场景演进}：v3是单星决策，v4扩展为多星协同。

\textbf{新增要素}：
\begin{itemize}
  \item \textbf{多卫星系统}：3颗遥感卫星（不同轨道）
  \item \textbf{协同覆盖}：某些大面积目标需要多星协同
  \item \textbf{资源竞争}：多星共享地面站下传资源
  \item \textbf{通信约束}：星间通信有延迟和带宽限制
\end{itemize}

\textbf{为什么v3方法不够？}
\begin{itemize}
  \item 单智能体规划无法处理智能体间的协调
  \item 需要考虑资源分配和任务分工
  \item 集中式规划可能不可行（通信限制）
\end{itemize}

\textbf{预习思考}：
\begin{enumerate}
  \item 3颗卫星如何分配10个目标？
  \item 如果某颗卫星观测失败，其他卫星能否补救？
  \item 集中式规划和分布式规划各有什么优缺点？
\end{enumerate}
\end{caseboxA}

% =============================================
\section{分布式规划（MA-STRIPS, 计划合并）}
% =============================================

\subsection{多智能体规划问题}

\begin{definition}[MA-STRIPS]
多智能体STRIPS问题扩展经典STRIPS：
\begin{itemize}
  \item 多个智能体 $\mathcal{AG} = \{ag_1, \ldots, ag_n\}$
  \item 每个智能体有自己的动作集合 $A_i$
  \item 某些谓词是公共的，某些是私有的
  \item 目标可能涉及多个智能体的状态
\end{itemize}
\end{definition}

\subsection{集中式 vs 分布式规划}

\begin{table}[htbp]
\centering
\caption{集中式与分布式规划对比}
\begin{tabular}{p{3cm}p{5cm}p{5cm}}
\toprule
\textbf{方面} & \textbf{集中式规划} & \textbf{分布式规划} \\
\midrule
决策方式 & 中央规划器统一决策 & 各智能体独立决策 \\
通信需求 & 需要全局信息 & 只需局部通信 \\
计算复杂度 & 指数增长于智能体数 & 可分解计算 \\
最优性 & 可保证全局最优 & 通常次优 \\
鲁棒性 & 单点故障风险 & 更鲁棒 \\
\bottomrule
\end{tabular}
\end{table}

\subsection{计划合并方法}

\textbf{计划合并}（Plan Merging）策略：
\begin{enumerate}[leftmargin=2em]
  \item 各智能体独立生成局部规划
  \item 检测规划间的冲突（资源冲突、目标冲突）
  \item 通过协商或约束传播解决冲突
  \item 合并为一致的联合规划
\end{enumerate}

% =============================================
\section{多智能体路径规划（MAPF, CBS）}
% =============================================

\subsection{MAPF问题定义}

\begin{definition}[多智能体路径规划（MAPF）]
在图 $G=(V,E)$ 上，$n$ 个智能体各有起点和终点，需要找到无冲突的路径使所有智能体从起点到达终点。

\textbf{冲突类型}：
\begin{itemize}
  \item \textbf{顶点冲突}：两个智能体同时占据同一顶点
  \item \textbf{边冲突}：两个智能体同时穿越同一条边（相向）
\end{itemize}
\end{definition}

\subsection{冲突搜索（CBS）}

\begin{keypoint}[CBS算法核心思想]
CBS（Conflict-Based Search）是求解最优MAPF的经典算法：
\begin{enumerate}
  \item \textbf{高层搜索}：在约束树中搜索，每个节点是一组约束集
  \item \textbf{低层搜索}：在给定约束下为每个智能体独立规划路径
  \item \textbf{冲突检测}：检查当前解中的冲突
  \item \textbf{分支}：对每个冲突生成两个子节点（分别约束两个智能体）
\end{enumerate}
\end{keypoint}

\begin{algorithm}[htbp]
\caption{CBS算法框架}
\begin{algorithmic}[1]
\Require MAPF问题实例
\Ensure 最优无冲突路径集
\State 初始化根节点：无约束，为每个智能体计算最短路径
\State 将根节点加入OPEN列表
\While{OPEN非空}
  \State 取出代价最小的节点 $N$
  \State 检测 $N$ 中的路径冲突
  \If{无冲突}
    \State \Return $N$ 中的路径集
  \EndIf
  \State 选择一个冲突 $(ag_i, ag_j, v, t)$
  \State 创建两个子节点：
  \State \quad $N_1$：增加约束"$ag_i$在时刻$t$不能在$v$"
  \State \quad $N_2$：增加约束"$ag_j$在时刻$t$不能在$v$"
  \State 为子节点重新规划受影响智能体的路径
  \State 将子节点加入OPEN
\EndWhile
\end{algorithmic}
\end{algorithm}

\subsection{优先级规划}

\textbf{优先级规划}是一种高效的次优方法：
\begin{enumerate}[leftmargin=2em]
  \item 为智能体分配优先级顺序
  \item 按优先级依次规划，高优先级智能体的路径作为低优先级的障碍
  \item 计算高效但不保证最优
\end{enumerate}

% =============================================
\section{\caseB{v4}多平台协同打击}
% =============================================

\begin{caseboxB}[案例B-v4：有人-无人协同作战]
\textbf{场景演进}：v3是单平台决策，v4扩展为多平台协同。

\textbf{新增要素}：
\begin{itemize}
  \item \textbf{异构平台}：1架有人机（指挥）+ 2架无人机（执行）
  \item \textbf{任务分工}：有人机负责决策，无人机负责侦察和打击
  \item \textbf{协同约束}：编队保持、通信链路、火力协调
\end{itemize}

\textbf{多智能体协调挑战}：
\begin{itemize}
  \item 如何分配侦察和打击任务给不同平台？
  \item 如何协调多个平台的时序配合？
  \item 通信中断时如何保持协同？
\end{itemize}
\end{caseboxB}

\subsection{任务分配问题}

多平台任务分配可建模为分配问题：

\begin{itemize}[leftmargin=2em]
  \item \textbf{集中式拍卖}：中央节点收集出价，分配任务
  \item \textbf{分布式拍卖}：智能体间直接协商
  \item \textbf{CBBA}（Consensus-Based Bundle Algorithm）：分布式任务分配算法
\end{itemize}

% =============================================
\section{多智能体强化学习（QMIX, MAPPO, CTDE）}
% =============================================

\subsection{MARL挑战}

多智能体强化学习面临特殊挑战：
\begin{itemize}[leftmargin=2em]
  \item \textbf{非平稳性}：其他智能体的策略在变化
  \item \textbf{信用分配}：团队奖励如何分配给个体？
  \item \textbf{可扩展性}：联合动作空间指数增长
  \item \textbf{通信}：是否需要智能体间通信？
\end{itemize}

\subsection{CTDE范式}

\begin{keypoint}[集中训练分布执行(CTDE)]
CTDE（Centralized Training Decentralized Execution）是MARL的主流范式：
\begin{itemize}
  \item \textbf{训练阶段}：可以访问全局信息（所有智能体的观测和动作）
  \item \textbf{执行阶段}：每个智能体只基于局部观测决策
  \item 兼顾学习效率和部署灵活性
\end{itemize}
\end{keypoint}

\subsection{QMIX}

QMIX是值分解方法的代表：

\begin{itemize}[leftmargin=2em]
  \item 每个智能体学习局部Q函数 $Q_i(o_i, a_i)$
  \item 全局Q函数是局部Q函数的单调混合：
  \[ Q_{\text{tot}} = f_{\text{mix}}(Q_1, \ldots, Q_n; s) \]
  \item 混合网络保证：$\frac{\partial Q_{\text{tot}}}{\partial Q_i} \geq 0$
  \item 执行时各智能体独立选择使 $Q_i$ 最大的动作
\end{itemize}

\subsection{MAPPO}

MAPPO将PPO扩展到多智能体设置：
\begin{itemize}[leftmargin=2em]
  \item 每个智能体有独立的策略网络
  \item 使用集中式Critic（输入全局状态）
  \item 在多个协作任务上表现优异
\end{itemize}

\subsection{案例应用：无人机蜂群}

\begin{appbox}[应用：无人机蜂群协同]
无人机蜂群任务规划是MARL的典型应用：
\begin{itemize}
  \item 多架无人机协同执行侦察、打击任务
  \item 需要处理编队控制、避障、任务分配
  \item MARL可以学习端到端的协同策略
  \item 国内研究：杜永浩等、贾高伟等的综述系统梳理了该领域进展
\end{itemize}
\end{appbox}

\subsection{HTN在多智能体协调中的应用}

华中科技大学王红卫教授、祁超教授团队将HTN规划思想应用于多智能体协调，在Applied Intelligence（2013）等期刊发表了资源增强HTN规划方法，提出了基于HTN的应急资源协调规划方法：

\begin{keypoint}[HTN协调规划的核心思想]
\begin{itemize}
  \item 将协调过程\textbf{嵌入}HTN规划过程，而非事后解决冲突
  \item 在任务分解时即考虑资源约束和智能体能力匹配
  \item 支持分布式规划环境下的实时资源冲突处理
\end{itemize}
这种方法体现了\textbf{模型}（HTN任务分解知识）与\textbf{算法}（协调机制）的深度融合。
\end{keypoint}

% =============================================
\section{本章小结与讨论}
% =============================================

\subsection*{本章小结}

\begin{itemize}[leftmargin=2em]
  \item 多智能体规划需要处理协调、通信和资源竞争
  \item MA-STRIPS扩展经典规划到多智能体设置
  \item CBS是求解最优MAPF的有效方法
  \item MARL通过学习获得协同策略，CTDE是主流范式
\end{itemize}

\begin{discussbox}
\textbf{讨论题}：
\begin{enumerate}
  \item 案例A-v4中，如果一颗卫星故障，系统应如何重新规划？
  \item CBS和优先级规划各适用于什么场景？
  \item 案例B-v4中的有人-无人协同，有人机和无人机的决策权限应如何分配？
  \item MARL学到的协同策略是否可解释？如何增强可解释性？
\end{enumerate}

\textbf{编程练习}：
\begin{enumerate}
  \item 实现CBS算法，在小规模MAPF实例上测试
  \item 使用PyMARL框架，在StarCraft微操环境中训练QMIX
  \item 将案例A-v4建模为多智能体任务分配问题并求解
\end{enumerate}
\end{discussbox}


% =============================================
% 第6章：博弈规划与对抗规划
% =============================================
% 第6章：博弈规划与对抗规划
\chapter{博弈规划与对抗规划}

\begin{levelbox}
本章介绍在对抗环境下的规划方法，包括博弈论基础、Stackelberg博弈、零和博弈以及欺骗与隐蔽规划。通过案例v5版本学习如何处理对抗性场景。建议学时：8学时。
\end{levelbox}

% =============================================
\section{博弈论基础}
% =============================================

\subsection{博弈的基本概念}

\begin{definition}[标准式博弈]
一个$n$人标准式博弈由以下要素定义：
\begin{itemize}
  \item 参与者集合 $N = \{1, 2, \ldots, n\}$
  \item 每个参与者的策略集合 $S_i$
  \item 每个参与者的收益函数 $u_i: S_1 \times \cdots \times S_n \rightarrow \mathbb{R}$
\end{itemize}
\end{definition}

\subsection{纳什均衡}

\begin{definition}[纳什均衡]
策略组合 $(s_1^*, \ldots, s_n^*)$ 是纳什均衡，当且仅当对每个参与者 $i$：
\[ u_i(s_i^*, s_{-i}^*) \geq u_i(s_i, s_{-i}^*), \quad \forall s_i \in S_i \]
即没有任何参与者有动机单方面偏离。
\end{definition}

\subsection{博弈类型}

\begin{table}[htbp]
\centering
\caption{常见博弈类型}
\begin{tabular}{p{3cm}p{4cm}p{5cm}}
\toprule
\textbf{类型} & \textbf{特点} & \textbf{应用场景} \\
\midrule
零和博弈 & 一方所得即另一方所失 & 军事对抗、竞技游戏 \\
Stackelberg博弈 & 领导者先行动，跟随者响应 & 安全资源部署、定价策略 \\
贝叶斯博弈 & 存在不完全信息 & 信息不对称的对抗 \\
重复博弈 & 同一博弈重复进行 & 长期关系、声誉建立 \\
\bottomrule
\end{tabular}
\end{table}

% =============================================
\section{Stackelberg博弈规划}
% =============================================

\subsection{Stackelberg博弈模型}

\begin{definition}[Stackelberg博弈]
Stackelberg博弈是一种序贯博弈：
\begin{enumerate}
  \item \textbf{领导者}先选择策略 $x$（可以是混合策略）
  \item \textbf{跟随者}观察到领导者策略后，选择最优响应 $y^*(x)$
  \item 领导者需要预测跟随者的响应来优化自己的策略
\end{enumerate}
\end{definition}

\textbf{Stackelberg均衡}：领导者选择使自己收益最大化的策略，同时考虑跟随者的最优响应。

\subsection{安全博弈}

\begin{keypoint}[安全博弈（Security Games）]
安全博弈是Stackelberg博弈的重要应用：
\begin{itemize}
  \item \textbf{防守方}（领导者）：部署有限的安全资源
  \item \textbf{攻击方}（跟随者）：观察防守部署后选择攻击目标
  \item 应用：机场安检、海岸巡逻、野生动物保护
\end{itemize}
\end{keypoint}

\begin{example}[案例A-v5的安全博弈建模]
\textbf{场景}：敌方可能对我方卫星实施电子干扰。

\textbf{防守方}（我方）：
\begin{itemize}
  \item 资源：有限的抗干扰设备和轨道机动能力
  \item 策略：为哪些卫星配置抗干扰？机动规避策略？
\end{itemize}

\textbf{攻击方}（敌方）：
\begin{itemize}
  \item 观察我方卫星轨道和配置
  \item 选择干扰目标和时机
\end{itemize}

\textbf{目标}：设计鲁棒的观测调度策略，即使面对最优攻击者也能保证最低观测收益。
\end{example}

% =============================================
\section{零和博弈与极大极小搜索}
% =============================================

\subsection{零和博弈}

\begin{definition}[零和博弈]
两人零和博弈中，双方收益之和恒为零：
\[ u_1(s_1, s_2) + u_2(s_1, s_2) = 0 \]
\end{definition}

零和博弈的最优策略可通过极大极小原则求解：
\[ v^* = \max_{s_1} \min_{s_2} u_1(s_1, s_2) = \min_{s_2} \max_{s_1} u_1(s_1, s_2) \]

\subsection{极大极小搜索}

\begin{algbox}[极大极小算法]
\begin{lstlisting}[language=Python,basicstyle=\ttfamily\small]
def minimax(state, depth, is_max_player):
    if depth == 0 or is_terminal(state):
        return evaluate(state)

    if is_max_player:
        value = -infinity
        for action in get_actions(state):
            child = apply(state, action)
            value = max(value, minimax(child, depth-1, False))
        return value
    else:
        value = +infinity
        for action in get_actions(state):
            child = apply(state, action)
            value = min(value, minimax(child, depth-1, True))
        return value
\end{lstlisting}
\end{algbox}

\subsection{Alpha-Beta剪枝}

Alpha-Beta剪枝通过维护搜索窗口 $[\alpha, \beta]$ 减少不必要的搜索：
\begin{itemize}[leftmargin=2em]
  \item $\alpha$：Max玩家已知的最优下界
  \item $\beta$：Min玩家已知的最优上界
  \item 当 $\alpha \geq \beta$ 时剪枝
\end{itemize}

% =============================================
\section{\caseA{v5}对抗性观测任务}
% =============================================

\begin{caseboxA}[案例A-v5：对抗性卫星观测]
\textbf{场景演进}：v4是合作性多星系统，v5引入对抗性敌方。

\textbf{对抗要素}：
\begin{itemize}
  \item \textbf{电子干扰}：敌方可干扰我方卫星的成像和通信
  \item \textbf{目标隐蔽}：敌方高价值目标可能采取伪装措施
  \item \textbf{轨道威胁}：敌方可能对我方卫星实施接近操作
\end{itemize}

\textbf{博弈建模}：
\begin{itemize}
  \item 我方选择观测调度策略（可以是随机策略）
  \item 敌方选择干扰策略
  \item 求解Stackelberg均衡或零和博弈均衡
\end{itemize}
\end{caseboxA}

% =============================================
\section{贝叶斯博弈与不完全信息}
% =============================================

\subsection{贝叶斯博弈}

\begin{definition}[贝叶斯博弈]
贝叶斯博弈在标准博弈基础上增加：
\begin{itemize}
  \item \textbf{类型空间} $\Theta_i$：参与者 $i$ 的可能类型
  \item \textbf{先验分布} $p(\theta)$：类型的先验概率分布
  \item 每个参与者知道自己的类型，但只知道其他参与者类型的分布
\end{itemize}
\end{definition}

\textbf{贝叶斯纳什均衡}：每个参与者在给定自己类型和对他人类型的信念下，选择最优策略。

\subsection{信息不对称下的规划}

在对抗性规划中，信息不对称是常态：
\begin{itemize}[leftmargin=2em]
  \item 我方不完全了解敌方的能力和意图（敌方类型）
  \item 需要基于敌方类型的先验分布进行规划
  \item 可以通过观测更新对敌方类型的信念
\end{itemize}

% =============================================
\section{\caseB{v5}欺骗与隐蔽规划}
% =============================================

\begin{caseboxB}[案例B-v5：欺骗与隐蔽战术]
\textbf{场景演进}：v4是对称信息下的协同，v5引入欺骗与隐蔽。

\textbf{新增要素}：
\begin{itemize}
  \item \textbf{欺骗机动}：通过虚假行动误导敌方
  \item \textbf{意图隐蔽}：隐藏真实攻击意图
  \item \textbf{佯攻}：分散敌方注意力
\end{itemize}

\textbf{核心问题}：
\begin{itemize}
  \item 如何设计欺骗策略使敌方形成错误判断？
  \item 如何在隐蔽行动的同时达成作战目标？
  \item 欺骗是否违反伦理？军事领域的欺骗合法性？
\end{itemize}
\end{caseboxB}

\subsection{欺骗规划}

\begin{definition}[欺骗规划]
\textbf{欺骗规划}（Deceptive Planning）旨在使对手对规划者的目标或能力形成错误信念。
\end{definition}

\textbf{欺骗策略类型}：
\begin{enumerate}[leftmargin=2em]
  \item \textbf{目标欺骗}：隐藏真实目标，展示虚假目标
  \item \textbf{能力欺骗}：隐藏或夸大自身能力
  \item \textbf{意图欺骗}：展示与真实意图不同的行动计划
\end{enumerate}

\begin{keypoint}[欺骗规划的形式化]
给定观察者模型，欺骗规划寻找：
\begin{itemize}
  \item 使观察者对真实目标 $G^*$ 的置信度降低
  \item 使观察者对虚假目标 $G'$ 的置信度升高
  \item 同时仍能达成真实目标 $G^*$
\end{itemize}
\end{keypoint}

\subsection{隐蔽规划}

\begin{definition}[隐蔽规划]
\textbf{隐蔽规划}（Covert Planning）旨在达成目标的同时，使对手难以推断规划者的真实目标。
\end{definition}

与欺骗规划不同，隐蔽规划不主动展示虚假信息，而是最小化信息泄露。

\textbf{隐蔽性度量}：
\begin{itemize}[leftmargin=2em]
  \item 目标概率分布的熵：熵越高，目标越难推断
  \item 观察者预测与真实目标的差异
\end{itemize}

\subsection{案例B-v5示例}

\begin{example}[佯攻与主攻]
\textbf{场景}：需要攻击高价值目标T1，但T1防御严密。

\textbf{欺骗策略}：
\begin{enumerate}
  \item 派遣无人机对次要目标T2实施明显的侦察行动
  \item 使敌方认为T2是主攻目标，调动防御资源
  \item 趁机对T1实施突然打击
\end{enumerate}

\textbf{隐蔽策略}：
\begin{enumerate}
  \item 对多个目标（T1、T2、T3）同时进行低强度侦察
  \item 使敌方无法判断真实攻击目标
  \item 在最后时刻才暴露真实意图
\end{enumerate}
\end{example}

% =============================================
\section{对抗多智能体系统}
% =============================================

\subsection{混合合作-竞争环境}

现实中很多场景既有合作又有竞争：
\begin{itemize}[leftmargin=2em]
  \item 队内合作、队间竞争（如足球比赛）
  \item 部分目标一致、部分目标冲突
  \item 可能存在背叛和策略转换
\end{itemize}

\subsection{对抗性MARL}

\begin{itemize}[leftmargin=2em]
  \item \textbf{自我对弈}（Self-Play）：通过与自己的历史版本对弈学习
  \item \textbf{虚拟自我对弈}（Fictitious Self-Play）：对历史策略分布采样
  \item \textbf{Policy-Space Response Oracles（PSRO）}：在策略空间中迭代求解
\end{itemize}

% =============================================
\section{本章小结与讨论}
% =============================================

\subsection*{本章小结}

\begin{itemize}[leftmargin=2em]
  \item 博弈论为对抗性规划提供数学框架
  \item Stackelberg博弈适用于领导者-跟随者场景
  \item 零和博弈通过极大极小搜索求解
  \item 欺骗规划和隐蔽规划是对抗性规划的重要分支
\end{itemize}

\begin{discussbox}
\textbf{讨论题}：
\begin{enumerate}
  \item 案例A-v5中，如果我方的观测调度是确定性的，敌方是否更容易干扰？随机策略有什么优势？
  \item 欺骗规划在民用领域有哪些应用？是否存在伦理问题？
  \item 案例B-v5的佯攻策略，如果敌方也是智能的会如何应对？会形成什么样的博弈？
  \item 如何评估欺骗规划的有效性？需要什么样的指标？
\end{enumerate}

\textbf{编程练习}：
\begin{enumerate}
  \item 实现极大极小搜索和Alpha-Beta剪枝，在井字棋上测试
  \item 使用GAMUT或Gambit工具求解案例A-v5的安全博弈
  \item 设计一个简单的欺骗规划算法，在路径规划场景中测试
\end{enumerate}
\end{discussbox}


% =============================================
% 第7章：前沿方法与系统集成
% =============================================
% 第7章：前沿方法与系统集成
\chapter{前沿方法与系统集成}

\begin{levelbox}
本章介绍任务规划的前沿方法（大语言模型、神经符号方法）和系统集成实践，通过综合案例展示完整的任务规划系统设计。建议学时：6学时。
\end{levelbox}

% =============================================
\section{大语言模型与规划}
% =============================================

\subsection{LLM的规划能力与局限}

大语言模型（LLM）在自然语言理解和生成方面表现出色，近年来也被用于任务规划。

\begin{keypoint}[LLM规划的优势与局限]
\textbf{优势}：
\begin{itemize}
  \item 自然语言接口，易于交互
  \item 丰富的世界知识
  \item 能够理解复杂的任务描述
  \item 具有一定的常识推理能力
\end{itemize}

\textbf{局限}：
\begin{itemize}
  \item 规划正确性难以保证（可能产生不可执行的计划）
  \item 复杂约束处理能力有限
  \item 长程规划容易出错
  \item 缺乏形式化验证能力
\end{itemize}
\end{keypoint}

\subsection{LLM+P：结合LLM与规划器}

\begin{definition}[LLM+P框架]
LLM+P利用LLM的语言理解能力和传统规划器的规划能力：
\begin{enumerate}
  \item LLM将自然语言任务描述转换为PDDL问题
  \item 调用Fast Downward等规划器求解
  \item LLM将规划结果转换回自然语言
\end{enumerate}
\end{definition}

\begin{lstlisting}[caption=LLM+P流程示例]
输入：自然语言任务描述
"卫星需要先观测北京区域，然后观测上海区域，最后返回待机状态"

LLM转换为PDDL：
(:goal (and
  (observed beijing)
  (observed shanghai)
  (standby satellite1)
  (before (observed beijing) (observed shanghai))
))

规划器求解，返回动作序列：
1. move(satellite1, pos1, beijing_orbit)
2. observe(satellite1, beijing)
3. move(satellite1, beijing_orbit, shanghai_orbit)
4. observe(satellite1, shanghai)
5. enter_standby(satellite1)

LLM转换回自然语言：
"首先将卫星移动到北京上空轨道位置，执行北京区域观测，
然后转移到上海上空，完成上海区域观测，最后进入待机模式。"
\end{lstlisting}

\subsection{思维链与推理增强}

\textbf{思维链}（Chain-of-Thought, CoT）：让LLM逐步推理，显式展示思考过程。

\textbf{思维树}（Tree of Thoughts, ToT）：扩展CoT，允许分支探索和回溯。

\textbf{ReAct}：结合推理（Reasoning）和行动（Acting），LLM可以调用外部工具获取信息。

\begin{example}[ReAct在任务规划中的应用]
\begin{verbatim}
思考：需要规划卫星观测任务，首先查询目标区域的云量
动作：调用气象API获取北京区域未来24小时云量预报
观测：北京区域预计14:00-16:00云量最低(20%)
思考：应该安排在14:00-16:00窗口观测北京
动作：检查卫星在该时段是否经过北京上空
观测：卫星在14:30-15:00经过北京上空
思考：时间窗口匹配，可以安排观测
动作：生成观测指令...
\end{verbatim}
\end{example}

\subsection{国内研究进展}

\begin{appbox}[title={\textbf{国内LLM规划研究}}]
\begin{itemize}
  \item \textbf{AutoPlan}（秦龙等, 自动化学报2024）：面向复杂任务的自主规划框架，结合LLM与规划引擎
  \item \textbf{MetaGPT}（深度求索DeepWisdom团队）：多智能体软件开发框架
  \item \textbf{ChatDev}（清华NLP实验室）：基于对话的软件开发智能体
  \item 国内大模型（文心一言、通义千问、ChatGLM、DeepSeek）在规划任务上持续优化
\end{itemize}
\end{appbox}

% =============================================
\section{神经符号规划与图神经网络}
% =============================================

\subsection{神经符号方法}

\begin{definition}[神经符号AI]
神经符号AI结合神经网络的学习能力和符号AI的推理能力：
\begin{itemize}
  \item 神经网络：感知、特征提取、模式识别
  \item 符号系统：逻辑推理、规划、可解释性
\end{itemize}
\end{definition}

\textbf{在规划中的应用}：
\begin{itemize}[leftmargin=2em]
  \item 神经网络学习启发式函数
  \item 神经网络生成动作候选
  \item 符号规划器保证解的正确性
\end{itemize}

\subsection{图神经网络在规划中的应用}

规划问题天然具有图结构（状态转移图、任务依赖图），图神经网络（GNN）可以有效捕获这些结构信息。

\textbf{应用场景}：
\begin{itemize}[leftmargin=2em]
  \item 学习PDDL问题的表示
  \item 预测动作的效用
  \item 生成规划启发式
\end{itemize}

% =============================================
\section{【综合案例】完整任务规划系统设计}
% =============================================

\begin{caseboxA}[title={\textbf{综合案例：卫星观测任务规划系统}}]
\textbf{系统目标}：设计一个完整的卫星观测任务规划系统，集成前6章的方法。

\textbf{系统需求}：
\begin{enumerate}
  \item 支持自然语言任务输入
  \item 处理多星协同调度
  \item 考虑不确定性（云层）
  \item 支持对抗性场景分析
  \item 提供可解释的规划结果
\end{enumerate}
\end{caseboxA}

\subsection{系统架构设计}

\begin{figure}[htbp]
\centering
\begin{tikzpicture}[
  box/.style={rectangle, draw, rounded corners, minimum width=3cm, minimum height=0.8cm, align=center},
  arrow/.style={->, thick}
]
  % 输入层
  \node[box, fill=blue!10] (input) at (0,4) {自然语言输入\\LLM解析};

  % 建模层
  \node[box, fill=green!10] (model) at (0,2.5) {PDDL建模\\约束生成};

  % 规划层
  \node[box, fill=yellow!10] (plan1) at (-3,1) {经典规划器\\(Fast Downward)};
  \node[box, fill=yellow!10] (plan2) at (0,1) {HTN规划器\\(SHOP2)};
  \node[box, fill=yellow!10] (plan3) at (3,1) {约束求解器\\(SMT)};

  % 优化层
  \node[box, fill=orange!10] (opt) at (0,-0.5) {多智能体协调\\博弈分析};

  % 输出层
  \node[box, fill=red!10] (output) at (0,-2) {规划结果\\可视化与解释};

  % 箭头
  \draw[arrow] (input) -- (model);
  \draw[arrow] (model) -- (plan1);
  \draw[arrow] (model) -- (plan2);
  \draw[arrow] (model) -- (plan3);
  \draw[arrow] (plan1) -- (opt);
  \draw[arrow] (plan2) -- (opt);
  \draw[arrow] (plan3) -- (opt);
  \draw[arrow] (opt) -- (output);
\end{tikzpicture}
\caption{卫星观测任务规划系统架构}
\end{figure}

\subsection{各模块说明}

\textbf{1. 输入解析模块}
\begin{itemize}[leftmargin=2em]
  \item 使用LLM理解自然语言任务描述
  \item 提取目标、约束、优先级等信息
  \item 生成结构化的任务规格
\end{itemize}

\textbf{2. 问题建模模块}
\begin{itemize}[leftmargin=2em]
  \item 根据任务规格生成PDDL/HDDL模型
  \item 编码时间窗口、资源约束
  \item 考虑不确定性因素的概率模型
\end{itemize}

\textbf{3. 规划求解模块}
\begin{itemize}[leftmargin=2em]
  \item 根据问题特性选择合适的规划器
  \item 简单问题：经典规划器
  \item 层次任务：HTN规划器
  \item 复杂约束：SMT求解器
\end{itemize}

\textbf{4. 多智能体协调模块}
\begin{itemize}[leftmargin=2em]
  \item 协调多星任务分配
  \item 解决资源冲突
  \item 对抗场景下的博弈分析
\end{itemize}

\textbf{5. 输出与解释模块}
\begin{itemize}[leftmargin=2em]
  \item 生成可执行的任务指令
  \item 提供规划结果的可视化
  \item 解释规划决策的理由
\end{itemize}

% =============================================
\section{可解释性与人机协同}
% =============================================

\subsection{可解释规划}

规划结果的可解释性对用户信任和人机协同至关重要。

\textbf{解释类型}：
\begin{enumerate}[leftmargin=2em]
  \item \textbf{为什么}：为什么选择这个动作？
  \item \textbf{为什么不}：为什么不选择另一个动作？
  \item \textbf{假设}：如果条件变化，规划会如何改变？
\end{enumerate}

\subsection{人机协同规划}

\begin{keypoint}[人机协同模式]
\begin{itemize}
  \item \textbf{人在环路}（Human-in-the-Loop）：人类审核和批准规划结果
  \item \textbf{混合主动性}：人和机器根据情况动态分配决策权
  \item \textbf{可调自主性}：根据任务关键程度调整自动化水平
\end{itemize}
\end{keypoint}

% =============================================
\section{规划系统的安全性}
% =============================================

\subsection{安全关键规划}

在安全关键应用中（航空、军事、医疗），规划系统必须满足严格的安全要求。

\textbf{安全性保障方法}：
\begin{enumerate}[leftmargin=2em]
  \item \textbf{形式化验证}：使用模型检测验证规划的正确性
  \item \textbf{安全约束}：硬编码安全边界，不允许违反
  \item \textbf{监控与回退}：运行时监控，异常时回退到安全状态
  \item \textbf{冗余设计}：多个独立规划器交叉验证
\end{enumerate}

\subsection{对抗鲁棒性}

在对抗环境中，需要考虑攻击者可能利用规划系统的漏洞。

\textbf{潜在攻击}：
\begin{itemize}[leftmargin=2em]
  \item 输入扰动：修改任务描述误导规划器
  \item 模型污染：污染学习数据影响学习型规划器
  \item 对抗样本：针对神经网络组件的对抗攻击
\end{itemize}

% =============================================
\section{本章小结与展望}
% =============================================

\subsection*{本章小结}

\begin{itemize}[leftmargin=2em]
  \item LLM为任务规划提供了自然语言接口
  \item LLM+P等混合方法结合了LLM和传统规划器的优势
  \item 神经符号方法是未来规划研究的重要方向
  \item 完整的规划系统需要集成多种方法
  \item 可解释性和安全性是实际部署的关键考量
\end{itemize}

\subsection*{未来展望}

\begin{enumerate}[leftmargin=2em]
  \item \textbf{更强的LLM规划能力}：随着LLM的发展，直接规划能力将增强
  \item \textbf{更紧密的神经-符号融合}：端到端可微的规划系统
  \item \textbf{更广泛的多模态规划}：结合视觉、语言、物理的规划
  \item \textbf{更可靠的安全保障}：形式化验证与AI的结合
\end{enumerate}

\begin{discussbox}
\textbf{讨论题}：
\begin{enumerate}
  \item LLM直接生成规划与LLM+P方法各有什么优缺点？未来会如何发展？
  \item 如果你要设计一个实际的卫星观测规划系统，会如何权衡自动化程度和人工干预？
  \item 规划系统的可解释性和性能之间是否存在权衡？如何平衡？
  \item 任务规划技术在军事应用中应该遵循什么伦理准则？
\end{enumerate}

\textbf{课程项目}：
\begin{enumerate}
  \item 选择案例A或案例B，设计并实现一个简化版的任务规划系统
  \item 系统应至少包含：问题建模、规划求解、结果展示三个模块
  \item 撰写项目报告，说明设计思路、实现方法和测试结果
\end{enumerate}
\end{discussbox}


% =============================================
% 附录
% =============================================
\appendix
% 附录A：算法分类速查表与学习路线图
\chapter{算法分类速查表与学习路线图}

本附录汇总全书涉及的主要算法，提供分类速查和学习建议。

% =============================================
\section*{A.1 算法总表}

\begin{longtable}{p{3.5cm}p{4cm}p{1.2cm}p{4.5cm}}
\toprule
\textbf{算法类别} & \textbf{代表算法} & \textbf{层次} & \textbf{核心特点} \\
\midrule
\endfirsthead
\toprule
\textbf{算法类别} & \textbf{代表算法} & \textbf{层次} & \textbf{核心特点} \\
\midrule
\endhead
\bottomrule
\endfoot

\multicolumn{4}{l}{\textbf{第一篇：经典规划方法}} \\
\midrule
无信息搜索 & BFS, DFS, IDS & \textcolor{ugcolor}{本} & 不使用启发信息 \\
启发式搜索 & A*, IDA*, GBFS & \textcolor{ugcolor}{本} & 启发函数引导搜索 \\
加权A* & WA* & \textcolor{gradcolor}{研} & 牺牲最优性换速度 \\
松弛启发式 & $h^{\max}$, $h^{\mathrm{add}}$, $h^{\mathrm{FF}}$ & \textcolor{ugcolor}{本} & 删除松弛估计 \\
规划图 & GraphPlan & \textcolor{ugcolor}{本} & 互斥分析+解提取 \\
抽象启发式 & PDB, Merge\&Shrink & \textcolor{gradcolor}{研} & 问题投影和抽象 \\
地标启发式 & LM-Cut, 地标计数 & \textcolor{gradcolor}{研} & 必经之路分析 \\
SAT规划 & SATplan, Madagascar & \textcolor{ugcolor}{本} & 转化为SAT求解 \\
CSP/SMT规划 & CSP编码, SMT规划 & \textcolor{ugcolor}{本}/\textcolor{gradcolor}{研} & 约束满足求解 \\
HTN规划 & SHOP/SHOP2, PANDA & \textcolor{ugcolor}{本} & 层次任务分解 \\
偏序规划 & POP & \textcolor{ugcolor}{本} & 最小承诺策略 \\
时态规划 & TFD, OPTIC & \textcolor{gradcolor}{研} & 处理持续性动作 \\
MDP/POMDP & 值迭代, PBVI & \textcolor{gradcolor}{研} & 概率不确定性 \\
FF规划器 & FF (EHC+$h^{\mathrm{FF}}$) & \textcolor{ugcolor}{本} & 松弛启发+爬山 \\
Fast Downward & FD (SAS$^+$) & \textcolor{ugcolor}{本} & 模块化框架 \\
LAMA & 地标+多启发式 & \textcolor{ugcolor}{本} & 满足规划冠军 \\
Scorpion/SymK & 饱和代价/符号搜索 & \textcolor{gradcolor}{研} & 最优规划前沿 \\
\midrule

\multicolumn{4}{l}{\textbf{第二篇：元启发式与优化方法}} \\
\midrule
遗传算法 & GA & \textcolor{ugcolor}{本} & 选择交叉变异 \\
模拟退火 & SA & \textcolor{ugcolor}{本} & 概率接受+冷却 \\
差分进化 & DE & \textcolor{gradcolor}{研} & 差异向量驱动 \\
禁忌搜索 & TS & \textcolor{gradcolor}{研} & 禁忌表+特赦 \\
粒子群优化 & PSO & \textcolor{gradcolor}{研} & 群体信息共享 \\
蚁群优化 & ACO, MMAS & \textcolor{gradcolor}{研} & 信息素引导 \\
新型元启发式 & GWO, WOA, HHO & \textcolor{selfcolor}{拓} & 自然现象启发 \\
\midrule

\multicolumn{4}{l}{\textbf{第三篇：强化学习与规划}} \\
\midrule
Dyna架构 & Dyna-Q & \textcolor{ugcolor}{本} & 学习+规划统一 \\
蒙特卡洛树搜索 & MCTS, UCT & \textcolor{ugcolor}{本} & 采样+模拟评估 \\
AlphaGo/Zero & MCTS+深度网络 & \textcolor{gradcolor}{研} & 自我对弈+搜索 \\
MuZero & 学习的世界模型 & \textcolor{gradcolor}{研} & 隐空间规划 \\
Dreamer系列 & DreamerV1--V3 & \textcolor{gradcolor}{研} & 想象中的规划 \\
DQN系列 & DQN, Rainbow & \textcolor{ugcolor}{本} & 深度Q网络 \\
策略梯度 & REINFORCE, A2C & \textcolor{ugcolor}{本} & 直接优化策略 \\
PPO & PPO & \textcolor{gradcolor}{研} & 裁剪策略更新 \\
SAC & SAC & \textcolor{gradcolor}{研} & 最大熵RL \\
层次RL & Options, HIRO & \textcolor{gradcolor}{研} & 时间抽象 \\
Decision Transformer & DT & \textcolor{gradcolor}{研} & 序列建模=RL \\
\midrule

\multicolumn{4}{l}{\textbf{第四篇：深度学习与规划}} \\
\midrule
图神经网络规划 & ASNets & \textcolor{gradcolor}{研} & 图结构泛化 \\
Transformer规划 & Plansformer & \textcolor{gradcolor}{研} & 端到端生成 \\
扩散模型规划 & Diffuser, Diff.\ Policy & \textcolor{selfcolor}{拓} & 轨迹去噪生成 \\
神经符号规划 & NLM, NSRT & \textcolor{selfcolor}{拓} & 感知+推理融合 \\
\midrule

\multicolumn{4}{l}{\textbf{第五篇：大语言模型与规划}} \\
\midrule
LLM直接规划 & Zero/Few-shot规划 & \textcolor{gradcolor}{研} & 语言知识驱动 \\
LLM+规划器 & LLM+P & \textcolor{gradcolor}{研} & 混合求解 \\
思维链/树/图 & CoT, ToT, GoT & \textcolor{gradcolor}{研} & 结构化推理 \\
具身LLM规划 & SayCan, Code as Policies & \textcolor{selfcolor}{拓} & 机器人+LLM \\
语言智能体 & ReAct, Reflexion & \textcolor{gradcolor}{研} & 推理+行动+反思 \\
多智能体LLM & AutoGen, MetaGPT & \textcolor{selfcolor}{拓} & 角色协作 \\
自主探索 & Voyager & \textcolor{selfcolor}{拓} & 终身学习 \\
\midrule

\multicolumn{4}{l}{\textbf{第六篇：多智能体规划}} \\
\midrule
MAPF & CBS, ICTS & \textcolor{ugcolor}{本}/\textcolor{gradcolor}{研} & 无冲突路径 \\
LaCAM & LaCAM/LaCAM* & \textcolor{gradcolor}{研} & 大规模MAPF \\
博弈论方法 & Nash, Stackelberg & \textcolor{gradcolor}{研} & 利益冲突建模 \\
MARL & QMIX, MAPPO & \textcolor{gradcolor}{研} & 集中训练分布执行 \\
分布式规划 & MA-STRIPS & \textcolor{selfcolor}{拓} & 隐私保护协调 \\
\end{longtable}

% =============================================
\section*{A.2 推荐学习路线}

\subsection*{本科生核心路线（建议32--48学时）}

\begin{enumerate}
  \item \textbf{基础（8--12学时）：}第1章（搜索基础）$\rightarrow$ 第2章1--2节（松弛启发式、规划图）
  \item \textbf{经典算法（8--12学时）：}第3章1--4节（SAT、HTN、POP）$\rightarrow$ 第4章1--3节（FF、FD、LAMA）
  \item \textbf{优化方法（4--6学时）：}第5章1--2节（GA、SA）
  \item \textbf{强化学习（8--12学时）：}第6章1--2节（Dyna、MCTS）$\rightarrow$ 第7章1--2节（DQN、策略梯度）
  \item \textbf{多智能体（4--6学时）：}第11章1节（MAPF、CBS）
\end{enumerate}

\subsection*{研究生扩展路线（在本科基础上增加16--32学时）}

\begin{enumerate}
  \item \textbf{高级启发式（4--6学时）：}第2章3--4节（PDB、LM-Cut）
  \item \textbf{时态与不确定性（4--6学时）：}第3章5--6节
  \item \textbf{高级RL（4--8学时）：}第6章3--4节（AlphaZero、Dreamer）$\rightarrow$ 第7章3--4节（层次RL、DT）
  \item \textbf{深度学习规划（4--6学时）：}第8章1--2节（GNN、Transformer）
  \item \textbf{LLM规划（4--8学时）：}第9章1--3节 $\rightarrow$ 第10章1--2节
  \item \textbf{多智能体进阶（4--6学时）：}第11章2--3节（博弈论、MARL）
\end{enumerate}

\subsection*{自学拓展专题}

以下内容适合感兴趣的读者根据研究方向选择阅读：
\begin{itemize}
  \item 扩散模型与规划（第8章3节）
  \item 神经符号规划（第8章4节）
  \item 具身智能LLM规划（第9章4节）
  \item 语言智能体前沿（第10章3--4节）
  \item 分布式规划（第11章4节）
  \item 新型元启发式综述（第5章4节）
\end{itemize}

% =============================================
\section*{A.3 算法演化脉络图}

任务规划算法的发展可以概括为以下几条主线：

\begin{enumerate}
  \item \textbf{符号推理主线：}STRIPS $\rightarrow$ GraphPlan $\rightarrow$ SAT规划 $\rightarrow$ 启发式搜索（FF/FD） $\rightarrow$ 地标/抽象启发式
  \item \textbf{层次分解主线：}HTN $\rightarrow$ SHOP/SHOP2 $\rightarrow$ PANDA $\rightarrow$ 时态HTN
  \item \textbf{优化求解主线：}SA/GA $\rightarrow$ PSO/ACO $\rightarrow$ 混合元启发式
  \item \textbf{学习规划主线：}Dyna $\rightarrow$ DQN $\rightarrow$ AlphaGo/MuZero $\rightarrow$ Dreamer $\rightarrow$ Decision Transformer
  \item \textbf{深度生成主线：}VAE/GAN $\rightarrow$ Transformer $\rightarrow$ 扩散模型 $\rightarrow$ 轨迹生成规划
  \item \textbf{语言智能主线：}GPT $\rightarrow$ CoT/ToT $\rightarrow$ ReAct/Reflexion $\rightarrow$ Voyager $\rightarrow$ 多智能体LLM
\end{enumerate}

这些主线并非孤立发展，而是在不断交叉融合。例如，MuZero结合了学习和搜索，LLM+P结合了语言模型和符号规划，ASNets结合了深度学习和经典规划。未来的发展趋势是各主线的进一步融合，形成更强大的统一规划框架。

% =============================================
\section*{A.4 国内推荐参考文献}

\subsection*{教材与专著}
\begin{enumerate}
  \item 蔡自兴, 刘丽珏, 蔡竞峰.《人工智能及其应用》（第7版）. 清华大学出版社, 2024. ——国内使用最广泛的AI教材之一，含智能规划章节。
  \item 吴飞, 潘云鹤.《人工智能引论》. 高等教育出版社, 2024. ——``101计划''核心教材，体现国内AI教育前沿。
  \item Ghallab等著, 殷建平等译.《自动规划：理论和实践》. 清华大学出版社, 2008. ——国际经典教材中文版，系统讲述自动规划理论。
  \item 段海滨.《蚁群算法原理及其应用》. 科学出版社, 2005. ——蚁群算法的理论与工程应用。
  \item 焦李成等.《免疫优化计算、学习与识别》. 科学出版社, 2006. ——免疫计算与优化理论。
  \item 焦李成等.《深度学习、优化与识别》. 清华大学出版社, 2017. ——深度学习与优化方法。
  \item 陈杰, 辛斌.《智能优化的探索--开发权衡理论与方法》. 科学出版社, 2013. ——元启发式核心理论。
  \item 陈杰, 方浩, 辛斌.《多智能体系统的协同群集运动控制》. 科学出版社, 2017. ——多智能体协同控制。
  \item 向峥嵘等.《多智能体协同控制及应用》. 化学工业出版社, 2023. ——多智能体系统理论与应用。
  \item 余伶俐, 周开军, 陈白帆.《智能驾驶技术：路径规划与导航控制》. 机械工业出版社, 2020. ——自动驾驶中的规划技术。
  \item 赵鑫等.《大语言模型》. 电子工业出版社, 2024. ——大语言模型技术全面参考。
\end{enumerate}

\subsection*{重要综述论文}
\begin{enumerate}
  \item 高阳, 陈世福, 陆鑫. 强化学习研究综述. 自动化学报, 2004. ——中文RL经典入门综述。
  \item 刘全等. 深度强化学习综述. 计算机学报, 2018. ——深度RL方法体系梳理。
  \item 赵冬斌等. 深度强化学习综述：兼论计算机围棋的发展. 控制理论与应用, 2016. ——AlphaGo技术解析。
  \item 梁星星, 冯旸赫等. 多Agent深度强化学习综述. 自动化学报, 2020. ——MARL方法全面综述。
  \item 杜永浩, 邢立宁, 蔡昭权. 无人飞行器集群智能调度技术综述. 自动化学报, 2020. ——无人机调度方法体系。
  \item 贾高伟, 王建峰. 无人机集群任务规划方法研究综述. 系统工程与电子技术, 2021. ——无人机集群规划综述。
  \item 邢立宁, 陈英武. 任务规划系统研究及发展. 火力与指挥控制, 2006. ——任务规划系统发展脉络。
  \item 秦龙等. 基于大语言模型的复杂任务自主规划处理框架. 自动化学报, 2024. ——LLM任务规划前沿。
  \item 张宏宏等. 有人/无人机协同作战：概念、技术与挑战. 航空学报, 2024. ——有人/无人协同综述。
  \item 王建峰等. 多无人机协同任务规划方法研究综述. 系统工程与电子技术, 2024. ——多无人机协同最新进展。
  \item 王雪松等. 安全强化学习综述. 自动化学报, 2023. ——安全RL约束方法。
\end{enumerate}


% ============ 参考文献 ============
\begin{thebibliography}{99}
\addcontentsline{toc}{chapter}{参考文献}

% 经典教材
\bibitem{ghallab2004} Ghallab M, Nau D, Traverso P. Automated Planning: Theory and Practice[M]. Morgan Kaufmann, 2004.
\bibitem{russell2020} Russell S, Norvig P. Artificial Intelligence: A Modern Approach (4th ed)[M]. Pearson, 2020.
\bibitem{sutton2018} Sutton R S, Barto A G. Reinforcement Learning: An Introduction (2nd ed)[M]. MIT Press, 2018.
\bibitem{shoham2009} Shoham Y, Leyton-Brown K. Multiagent Systems[M]. Cambridge University Press, 2009.

% 经典规划
\bibitem{hoffmann2001} Hoffmann J, Nebel B. The FF Planning System[J]. JAIR, 2001, 14:253--302.
\bibitem{helmert2006} Helmert M. The Fast Downward Planning System[J]. JAIR, 2006, 26:191--246.
\bibitem{erol1994} Erol K, Hendler J, Nau D S. HTN Planning: Complexity and Expressivity[C]. AAAI, 1994.

% 强化学习与规划
\bibitem{silver2017} Silver D, et al. Mastering the Game of Go without Human Knowledge[J]. Nature, 2017.
\bibitem{schrittwieser2020} Schrittwieser J, et al. Mastering Atari, Go, Chess and Shogi by Planning with a Learned Model[J]. Nature, 2020.
\bibitem{hafner2023} Hafner D, et al. Mastering Diverse Domains through World Models[J]. arXiv:2301.04104, 2023.

% LLM与规划
\bibitem{yao2023} Yao S, et al. ReAct: Synergizing Reasoning and Acting in Language Models[C]. ICLR, 2023.
\bibitem{liu2023llmp} Liu B, et al. LLM+P: Empowering Large Language Models with Optimal Planning Proficiency[J]. arXiv:2304.11477, 2023.

% 多智能体
\bibitem{sharon2015} Sharon G, et al. Conflict-Based Search for Optimal Multi-Agent Pathfinding[J]. AI, 2015.
\bibitem{rashid2018} Rashid T, et al. QMIX: Monotonic Value Function Factorisation for Deep Multi-Agent RL[C]. ICML, 2018.

% 博弈规划
\bibitem{tambe2011} Tambe M. Security and Game Theory: Algorithms, Deployed Systems, Lessons Learned[M]. Cambridge, 2011.
\bibitem{basilico2017} Basilico N, et al. Adversarial Patrolling Games[J]. AI, 2017.

% 中文文献
\bibitem{cai2024} 蔡自兴, 等. 人工智能及其应用（第7版）[M]. 清华大学出版社, 2024.
\bibitem{yin2008} Ghallab等著, 殷建平等译. 自动规划：理论和实践[M]. 清华大学出版社, 2008.
\bibitem{qin2024} 秦龙, 等. 基于大语言模型的复杂任务自主规划处理框架[J]. 自动化学报, 2024.
\bibitem{liang2020} 梁星星, 等. 多Agent深度强化学习综述[J]. 自动化学报, 2020.
\bibitem{du2020} 杜永浩, 邢立宁, 等. 无人飞行器集群智能调度技术综述[J]. 自动化学报, 2020.

\end{thebibliography}

\end{document}
